\chapter{طراحی و  شبیه سازی سامانه پیشنهادی}
\section{مقدمه}
\label{sec:chapter4_intro}

در فصل سوم، مبانی نظری یادگیری تقویتی عمیق، فرمول‌بندی ناوبری در قالب فرایند تصمیم‌گیری مارکوف و الگوریتم \lr{SAC} تشریح شد. در این فصل، این مبانی به چارچوب پیشنهادی ناوبری خودران اسکوتر چهارچرخ در محیط‌های دینامیکی تبدیل می‌شوند؛ چارچوبی که ساختار محیط یادگیری، ادراک، کنترل، فضای مشاهده و عمل، تابع پاداش، منطق ایمنی، معماری عامل و فرایند آموزش را به‌صورت یکپارچه تعریف می‌کند.

ابتدا معماری ادراک--تصمیم--کنترل و بستر \lr{ROS2}/\lr{Gazebo} معرفی می‌شود. سپس مدل اسکوتر، حسگرها، موانع دینامیکی و لایه کنترل حرکتی تشریح می‌گردند. در ادامه، زنجیره ادراک شامل بینایی استریو، آشکارسازی موانع، تخمین موقعیت و سرعت، تشخیص هدف و حسگرهای فراصوت بررسی می‌شود. در نهایت، فضای مشاهده و عمل، تابع پاداش، شرایط پایان اپیزود، سپر ایمنی، معماری «کنشگر--منتقد نرم و حافظه بلند کوتاه‌مدت%
\LTRfootnote{SAC-LSTM}» %
و آموزش معلم--دانشجو و مرحله‌ای ارائه می‌شوند تا مبنای ارزیابی فصل‌های بعد فراهم شود.


\section{معماری کلی سامانه ناوبری پیشنهادی}
\label{sec:proposed_navigation_architecture}

سامانه پیشنهادی دارای معماری ماژولار و لایه‌ای است تا اجزای آن مستقل توسعه و ارزیابی شوند و در عین حال یک حلقه تصمیم‌گیری پیوسته تشکیل دهند. این معماری سه بخش اصلی دارد:
\begin{enumerate}
	\item لایه ادراک و استخراج ویژگی؛
	\item لایه تصمیم‌گیری مبتنی بر یادگیری تقویتی؛
	\item لایه کنترل و اعمال فرمان حرکتی.
\end{enumerate}
لایه ادراک داده‌های حسگری را به ویژگی‌های مورد نیاز عامل تبدیل می‌کند؛ سیاست یادگرفته‌شده بر پایه مشاهده جاری کنش را می‌سازد؛ و لایه کنترل، کنش را با رعایت حدود حرکتی و قیود ایمنی به فرمان قابل اعمال به اسکوتر تبدیل می‌کند.

\subsection{نمای کلی حلقه ادراک--تصمیم--کنترل}
\label{subsec:perception_decision_control_loop}

هسته سامانه، حلقه ادراک--تصمیم--کنترل سازگار با فرایند تصمیم‌گیری مارکوف است. با توجه به دسترسی‌نداشتن عامل به حالت کامل محیط، تصمیم‌ها از مشاهده‌های حسگری و ادراکی استخراج می‌شوند.

در گام $t$، داده‌های وضعیت ربات، هدف، موانع و حسگرهای نزدیک‌برد در مشاهده $o_t$ سازمان‌دهی و به سیاست داده می‌شوند. ساختار \lr{LSTM} نیز اطلاعات نهفته در توالی ورودی‌های مربوط به موانع را برای تشخیص روند حرکت موانع به کار می‌گیرد. سیاست کنش پیوسته $a_t$، شامل سرعت خطی و زاویه‌ای مطلوب، را تولید می‌کند؛ کنش پس از بررسی حدود مجاز و قیود ایمنی اعمال می‌شود و محیط، مشاهده $o_{t+1}$ و پاداش $r_t$ را بر اساس پیشرفت، ایمنی و کیفیت حرکت بازمی‌گرداند. چرخه تا رسیدن به هدف، برخورد یا پایان زمان ادامه می‌یابد. رابطه کلی این فرایند در شکل \ref{fig:state_representation} نشان داده شده است.

\begin{figure}[ht]
	\centering
	\includegraphics[width=0.6\textwidth]{state_representation}
	
	\caption{چرخه کلی ادراک، تصمیم‌گیری و اعمال فرمان در سامانه ناوبری پیشنهادی}
	\label{fig:state_representation}
	
	\vspace{0.4em}
	
	{\small
		\href{https://mohammadabedi1179.github.io/thesis-archive/chapters/ch4/state_representation.html}
		{\textbf{مشاهده نسخه آنلاین شکل}}
		\qquad
		\qrcode[height=1.2cm]{https://mohammadabedi1179.github.io/thesis-archive/chapters/ch4/state_representation.html}
	}
\end{figure}

سیاست با بیشینه‌سازی امید ریاضی مجموع پاداش‌های تنزیل‌شده در رابطه \eqref{eq:cumulative_reward}، توالی تصمیم‌هایی را می‌آموزد که در بلندمدت حرکت ایمن، پایدار و مؤثر به سمت هدف ایجاد کنند.

\subsection{جایگاه \lr{ROS2} و \lr{Gazebo} در معماری شبیه‌سازی}
\label{subsec:ros2_gazebo_architecture}

بستر شبیه‌سازی بر پایه \lr{Gazebo} و \lr{ROS2} ساخته شده است. \lr{Gazebo} دینامیک اسکوتر، تعامل فیزیکی، موانع و حسگرها را شبیه‌سازی می‌کند و \lr{ROS2} تبادل داده میان ماژول‌های ادراک، تشکیل مشاهده، تصمیم‌گیری، کنترل و ارزیابی را بر عهده دارد. بدین‌ترتیب، داده حسگری به مشاهده تبدیل می‌شود، سیاست فرمان را تولید می‌کند و فرمان دوباره به مدل حرکتی اعمال می‌گردد.
\\

اجزای اصلی این بستر عبارت‌اند از:
\begin{itemize}
	\item \textbf{محیط شبیه‌سازی:} مدل اسکوتر، محیط، هدف، موانع دینامیکی و تعاملات فیزیکی؛
	\item \textbf{لایه حسگری و ادراک:} دریافت داده و استخراج موقعیت نسبی هدف، موانع و فاصله‌های نزدیک‌برد؛
	\item \textbf{لایه تشکیل مشاهده:} تبدیل خروجی ادراک به ساختار ثابت مناسب سیاست؛
	\item \textbf{لایه تصمیم‌گیری:} سیاست \lr{SAC-LSTM} برای تولید سرعت خطی و زاویه‌ای؛
	\item \textbf{لایه کنترل و اعمال فرمان:} تبدیل و محدودسازی فرمان پیش از اعمال به مدل حرکتی؛
	\item \textbf{لایه ارزیابی و ثبت داده:} ثبت پاداش، وضعیت اپیزود، برخورد، موفقیت و سایر شاخص‌ها.
\end{itemize}
این تفکیک، اصلاح مستقل ماژول‌ها و آموزش و ارزیابی یکپارچه عامل را ممکن می‌سازد.


%==========================================================
\section{محیط شبیه‌سازی و مدل‌سازی سامانه}
\label{sec:simulation_environment_modeling}

در این پژوهش، به‌منظور ایجاد بستری واقع‌گرایانه برای آموزش و ارزیابی سیاست ناوبری، مدل اسکوتر چهارچرخ محرک و اجزای اصلی سامانه در محیط شبیه‌ساز \lr{Gazebo} تعریف شده‌اند. مدل ربات با استفاده از فایل‌های \lr{URDF/Xacro} توسعه یافته و شامل توصیف اجزای مکانیکی، مشخصات جرمی و اینرسی، جای‌نشانی حسگرها، مدل برخورد، مدل ظاهری و رابط کنترلی چرخ‌ها است. هدف از این مدل‌سازی، فراهم‌کردن محیطی است که در آن عامل یادگیری تقویتی بتواند در یک حلقه بسته با ربات و محیط تعامل داشته باشد و اثر تصمیم‌های حرکتی خود را در قالب مشاهده بعدی و پاداش دریافت کند.

در این بخش، ابتدا نحوه مدل‌سازی محیط پویا و حرکت موانع دینامیکی تشریح می‌شود و سپس مدل فیزیکی ربات و چارچوب‌های مرجع آن مورد بررسی قرار می‌گیرد. در ادامه، اجزای اصلی شامل بدنه، چرخ‌ها، دسته، دوربین استریو و حسگرهای فراصوت معرفی می‌شوند. جزئیات مربوط به ارتباط نرم‌افزاری گره‌ها، موضوعات ارتباطی و اجرای سناریوهای آموزشی در فصل بعد ارائه خواهد شد.

همچنین لازم به ذکر است، برای تولید حرکت موانع دینامیکی، یک مدل حرکت تصادفی مبتنی بر نیروهای اجتماعی اصلاح‌شده پیاده‌سازی شد. در این مدل، هر مانع تمایل دارد با یک سرعت مطلوب در جهتی که به‌صورت تصادفی و پیوسته تغییر می‌کند حرکت کند. علاوه بر این، نیروهای دافعه نمایی میان مانع متحرک و ربات، سایر موانع، نشانگر هدف و مرز محیط اعمال می‌شوند. به‌منظور افزایش پایداری عددی و جلوگیری قطعی از برخورد، محدودیت شتاب و سرعت، تغییر جهت تصادفی در نزدیکی مرز و یک لایه جداسازی سخت نیز به مدل افزوده شده است. ازاین‌رو، مدل استفاده‌شده نسخه‌ای اصلاح‌شده و توسعه‌یافته بر مبنای مفهوم نیروهای اجتماعی است و با مدل کلاسیک نیروی اجتماعی تفاوت‌هایی دارد.


%----------------------------------------------------------
\subsection{مدل ربات اسکوتر}
\label{sec:robot_model}

مدل اسکوتر چهارچرخ محرک در این پژوهش با استفاده از توصیف‌گر \lr{URDF/Xacro} تعریف شده است. هدف از این مدل‌سازی، ایجاد یک نمایش هندسی و دینامیکی قابل استفاده در شبیه‌ساز \lr{Gazebo} است؛ به‌گونه‌ای که تعامل فیزیکی چرخ‌ها با سطح زمین، جای‌نشانی حسگرها و تولید داده‌های ادراکی مورد نیاز عامل یادگیری تقویتی فراهم شود. در ساختار مدل، هسته اصلی ربات در یک فایل مجزا تعریف شده و حسگرها به‌صورت ماژول‌های مستقل به مدل افزوده شده‌اند. این رویکرد باعث می‌شود تغییر یا توسعه هر بخش، مانند اصلاح مدل حسگر یا تغییر موقعیت نصب آن، بدون نیاز به بازنویسی کامل مدل ربات امکان‌پذیر باشد.

\begin{figure}[ht]
	\centerline{\includegraphics[width=10cm, height=5cm]{Robot_model}}
	\caption{نمای کلی اسکوتر مدل‌شده در شبیه‌ساز \lr{Gazebo}}
	\label{fig:Robot_model}
\end{figure}

%----------------------------------------------------------
\subsubsection{چارچوب‌های مرجع و ساختار کلی لینک‌ها}
\label{subsubsec:robot_reference_frames}

در مدل حاضر، لینک پایه (\lr{base\_link}) به‌عنوان چارچوب مرجع اصلی ربات در نظر گرفته شده است. مطابق مدل تعریف‌شده، این چارچوب در مرکز هندسی چهار چرخ قرار دارد. انتخاب این نقطه به‌عنوان مرجع اصلی باعث می‌شود موقعیت چرخ‌ها نسبت به مرکز ربات به‌صورت متقارن تعریف شود و تبدیل فرمان‌های حرکتی به سرعت چرخ‌ها با سادگی و پایداری بیشتری انجام گیرد.

بدنه اصلی ربات، با نام \lr{chassis}، از طریق یک مفصل ثابت به لنک پایه متصل شده است. برای حفظ موقعیت هندسی بدنه نسبت به مدل اولیه و در عین حال نگه‌داشتن لنک پایه در مرکز چرخ‌ها، یک جابه‌جایی ثابت میان این دو لینک اعمال شده است:

\begin{equation}
	{}^{base}\mathbf{T}_{chassis}:
	\quad
	\Delta x=-0.4,\qquad
	\Delta y=0,\qquad
	\Delta z=0.
	\label{eq:base_to_chassis_transform}
\end{equation}

همانطور که در شکل \label{fig:wheel_links} نشان داده شده است، اجزای حرکتی اصلی، یعنی چرخ‌ها، حول یک مرجع مرکزی تعریف می‌شوند، در حالی که بدنه، اجزای بصری و حجم برخورد در موقعیت مناسب خود نسبت به کل ربات باقی می‌مانند.

\begin{figure}[ht]
	\centering
	\subfloat[]{
		\includegraphics[width=0.45\textwidth]{wheel_links_down_view}
	}
	\hfill
	\subfloat[]{
		\includegraphics[width=0.45\textwidth]{wheel_links_3_corners_view}
	}
	\caption{نمای لینک‌های تعریف‌شده برای چرخ‌ها و بدنه: الف) نمای پایین، ب) نمای سه‌بعدی از اتصال چرخ‌ها به بدنه}
	\label{fig:wheel_links}
	
	\vspace{0.4em}
	
	{\small
		\href{https://mohammadabedi1179.github.io/thesis-archive/chapters/ch4/wheel_links_3_corners_view.html}
		{\textbf{مشاهده نسخه آنلاین شکل}}
		\qquad
		\qrcode[height=1.2cm]{https://mohammadabedi1179.github.io/thesis-archive/chapters/ch4/wheel_links_3_corners_view.html}
	}
\end{figure}

\begin{figure}[ht]

\end{figure}

%----------------------------------------------------------
\subsubsection{مدل‌سازی بدنه و پارامترهای دینامیکی}
\label{subsubsec:scooter_spec}

بدنه ربات با لینک شاسی تعریف شده است. برای مدل‌سازی جرم و اینرسی، از ابعاد اصلی بدنه و برآورد جرم مجموعه شامل شاسی، درایورها، باتری‌ها، بردهای الکترونیکی و حسگرها استفاده شده است. در مدل \lr{URDF/Xacro}، اینرسی بدنه با تقریب جعبه‌ای محاسبه شده و مرکز جرم آن نسبت به چارچوب شاسی با یک جابه‌جایی مشخص تعریف شده است. این تقریب باعث کاهش پیچیدگی مدل دینامیکی می‌شود و در عین حال برای هدف این پژوهش، یعنی آموزش و ارزیابی سیاست ناوبری، دقت کافی فراهم می‌کند.

\begin{table}[ht]
	\caption{مشخصات کلی بدنه اسکوتر}
	\label{tab:my-scooter-specs}
	\centering
	\onehalfspacing
	\begin{tabular}{c|c}
		\hline
		\multicolumn{1}{|c|}{کمیت}  & \multicolumn{1}{c|}{اندازه}   \\ \hline
		\multicolumn{1}{|c|}{جرم}   & \multicolumn{1}{c|}{$25~kg$}    \\ \hline
		\multicolumn{1}{|c|}{طول}   & \multicolumn{1}{c|}{$0.8~m$}    \\ \hline
		\multicolumn{1}{|c|}{عرض}   & \multicolumn{1}{c|}{$0.6108~m$} \\ \hline
		\multicolumn{1}{|c|}{ارتفاع بدنه}   & \multicolumn{1}{c|}{$0.15~m$} \\ \hline
		\multicolumn{1}{|c|}{ارتفاع با احتساب میله فرمان}   & \multicolumn{1}{c|}{$0.55~m$} \\ \hline
	\end{tabular}
\end{table}

برای کاهش هزینه محاسباتی برخورد در شبیه‌سازی، حجم «برخورد%
\LTRfootnote{collision}» %
 بدنه به‌صورت یک جعبه هم‌ابعاد با ابعاد اصلی شاسی مدل شده است. در مقابل، برای نمایش ظاهری از فایل مش استفاده شده تا شکل ربات از نظر بصری به مدل واقعی نزدیک‌تر باشد. این تفکیک میان مدل بصری و مدل برخورد، یک راهکار رایج در شبیه‌سازی رباتیکی است؛ زیرا ضمن حفظ نمایش واقع‌گرایانه، از افزایش غیرضروری هزینه محاسباتی در حل تماس‌ها جلوگیری می‌کند.

%----------------------------------------------------------
\subsubsection{مدل چرخ‌ها، مفاصل و جای‌گذاری هندسی}
\label{subsubsec:wheel_model_geometry}

ربات دارای چهار چرخ محرک با نام‌های \lr{FL\_link}، \lr{FR\_link}، \lr{BL\_link} و \lr{BR\_link} است که هر یک از طریق یک مفصل دورانی از نوع پیوسته به لینک پایه متصل شده‌اند. محور دوران چرخ‌ها در راستای محور \(y\)، یعنی \([0\;\;1\;\;0]^T\)، تعریف شده است. موقعیت چرخ‌ها نسبت به لینک پایه به‌صورت متقارن در جلو/عقب و چپ/راست مدل شده است.

\begin{table}[ht]
	\centering
	\caption{موقعیت نسبی چرخ‌های اسکوتر نسبت به لینک پایه}
	\label{tab:wheel_relative_positions}
	\begin{tabular}{|c|c|}
		\hline
		\textbf{نام چرخ} & \textbf{موقعیت نسبی \((x,y,z)\) بر حسب متر} \\
		\hline
		\text{جلو-چپ} & $(+0.3295, +0.33054, 0)$ \\
		\text{جلو-راست} & $(+0.3295, -0.33054, 0)$ \\
		\text{عقب-چپ} & $(-0.3295, +0.33054, 0)$ \\
		\text{عقب-راست} & $(-0.3295, -0.33054, 0)$ \\
		\hline
	\end{tabular}
\end{table}

برای هر چرخ، نمایش ظاهری با استفاده از فایل مش انجام شده است، در حالی که حجم برخورد چرخ‌ها با زمین به‌صورت یک استوانه ساده مدل شده است. شعاع چرخ‌ها برابر با \(r=0.0925~m\) و ضخامت مدل برخورد استوانه‌ای برابر با \(0.04~m\) در نظر گرفته شده است. این ساده‌سازی باعث می‌شود تماس چرخ با زمین در \lr{Gazebo} پایدارتر و محاسبات برخورد سبک‌تر شود.

در مدل هندسی اولیه چرخ‌ها، ناهمواری‌های سطحی برای افزایش اصطکاک با زمین وجود داشت. با این حال، استفاده مستقیم از چنین هندسه‌ای در محاسبات برخورد شبیه‌ساز، به‌ویژه در آموزش‌های طولانی‌مدت یادگیری تقویتی، باعث افزایش هزینه محاسباتی و کاهش پایداری عددی می‌شود. به همین دلیل، در مدل نهایی، هندسه برخورد چرخ‌ها به یک استوانه صاف ساده‌سازی شده و اثر تماس با زمین از طریق تنظیم ضرایب اصطکاک در شبیه‌ساز لحاظ شده است.

در فایل مدل، ضرایب اصطکاک تماس چرخ با زمین در دو راستای تماس به‌صورت \(\mu_1=1.0\) و \(\mu_2=0.1\) تعریف شده‌اند. انتخاب این ضرایب با هدف ایجاد کشش کافی در راستای اصلی حرکت و در عین حال حفظ پایداری عددی تماس در شبیه‌سازی انجام شده است. با توجه به حساسیت آموزش یادگیری تقویتی به ناپایداری‌های فیزیکی شبیه‌ساز، این ساده‌سازی تعادلی میان واقع‌گرایی، پایداری و سرعت آموزش ایجاد می‌کند.

%----------------------------------------------------------
\subsubsection{مدل دسته و ساختار نصب حسگرها}
\label{subsubsec:handle_sensor_mount}

برای نصب حسگرهای بالاتنه، به‌ویژه دوربین استریو، یک لینک میانی در نظر گرفته شده است. این لینک به‌صورت یک مفصل ثابت به لینک پایه متصل شده و در موقعیت \((0.3355,0,0.43)\) نسبت به چارچوب پایه قرار گرفته است. لینک مذکور به‌صورت یک استوانه با طول \(0.7~m\) و شعاع \(0.015~m\) مدل شده و نقش سازه نگهدارنده حسگرهای دیداری را ایفا می‌کند.

%----------------------------------------------------------
\subsubsection{حسگر دوربین استریو}
\label{subsubsec:stereo_camera_model}

دوربین استریو به‌صورت یک ماژول مستقل تعریف شده و به لینک میانی متصل شده است. این ماژول شامل یک لینک پشتیبان و دو دوربین چپ و راست است. لینک پشتیبان دوربین با یک مفصل ثابت به دسته متصل شده و دوربین‌های چپ و راست روی آن در راستای محور \(y\) قرار گرفته‌اند. مطابق مدل تعریف‌شده، موقعیت دوربین چپ و راست نسبت به لینک پشتیبان به‌ترتیب برابر با \(y=+0.03~m\) و \(y=-0.03~m\) است. بنابراین، فاصله مبنای هندسی دوربین‌ها برابر است با:

\begin{equation}
	b = 0.03 - (-0.03) = 0.06~m 
	\label{eq:stereo_camera_baseline}
\end{equation}

مشخصات شبیه‌سازی هر دوربین در \lr{Gazebo} شامل رزولوشن \(1640\times1232\)، نرخ به‌روزرسانی \(15~Hz\)، زاویه دید افقی \(1.0448~rad\) و بازه دید از \(0.02\) تا \(15~m\) است. خروجی دوربین‌ها در قالب موضوعات تصویری منتشر می‌شود و در زنجیره ادراک برای استخراج اطلاعات عمق و ویژگی‌های مرتبط با هدف و موانع مورد استفاده قرار می‌گیرد.

%----------------------------------------------------------
\subsubsection{حسگرهای فراصوت \lr{SRF05}}
\label{subsubsec:srf05_sensor_model}

برای افزایش ایمنی در فاصله نزدیک و پوشش نواحی‌ای که ممکن است توسط حسگر دیداری به‌خوبی پوشش داده نشوند، پنج حسگر نزدیک‌برد از نوع \lr{SRF05} در مدل ربات تعریف شده‌اند. این حسگرها شامل دو حسگر در بخش جلویی ربات، یک حسگر در سمت چپ، یک حسگر در سمت راست و یک حسگر در بخش عقب هستند. حسگرهای فراصوت به‌صورت لینک‌های مستقل به لینک پایه متصل شده‌اند و موقعیت آن‌ها نسبت به چارچوب مرکزی ربات در جدول \ref{tab:srf05_relative_positions} آمده است.

\begin{table}[ht]
	\centering
	\caption{موقعیت نسبی حسگرهای فراصوت نسبت به لینک پایه}
	\label{tab:srf05_relative_positions}
	\begin{tabular}{|c|c|}
		\hline
		\textbf{نام موقعیت} & \textbf{موقعیت نسبی \((x,y,z)\) بر حسب متر} \\
		\hline
		\text{جلو-چپ} & $(+0.275, +0.06, 0.01)$ \\
		\text{جلو-راست} & $(+0.275, -0.06, 0.01)$ \\
		\text{سمت چپ} & $(0.0, +0.175, 0.01)$ \\
		\text{سمت راست} & $(0.0, -0.175, 0.01)$ \\
		\text{عقب} & $(-0.275, 0.0, 0.01)$ \\
		\hline
	\end{tabular}
\end{table}

هر حسگر در \lr{Gazebo} به‌صورت یک حسگر \lr{ray} مدل شده است. نرخ تولید داده برای این حسگرها برابر با \(20~Hz\) در نظر گرفته شده و هر حسگر در بازه زاویه‌ای \(\pm 0.13~rad\) با تعداد ۵ نمونه در هر اسکن، فاصله موانع را اندازه‌گیری می‌کند. بازه اندازه‌گیری حسگرها از \(0.04\) تا \(3.0~m\) تعریف شده و تفکیک‌پذیری فاصله برابر با \(0.01~m\) است.

خروجی این حسگرها در بردار مشاهده عامل یادگیری تقویتی مورد استفاده قرار می‌گیرد. نقش اصلی این حسگرها، فراهم‌کردن اطلاعات نزدیک‌برد برای کاهش احتمال برخورد، تقویت تصمیم‌گیری در فاصله‌های بحرانی و افزایش پایداری رفتار عامل در شرایطی است که داده‌های دیداری به‌تنهایی برای تصمیم‌گیری ایمن کافی نیستند.
\section{لایه کنترل حرکتی و اعمال فرمان سرعت به چرخ‌ها}
\label{sec:velocity_control_layer}

در سامانه پیشنهادی، خروجی عامل یادگیری تقویتی به‌صورت فرمان سرعت خطی و سرعت زاویه‌ای ربات تعریف می‌شود. از آنجا که اسکوتر مورد بررسی دارای چهار چرخ محرک با محورهای ثابت است، فرمان سطح بالای عامل باید به سرعت زاویه‌ای هر یک از چرخ‌ها تبدیل شود. این تبدیل بر پایه سازوکار \lr{skid-steering} انجام می‌شود؛ به این معنا که اختلاف سرعت میان چرخ‌های سمت چپ و راست موجب ایجاد چرخش حول محور عمودی ربات می‌گردد.

در مدل شبیه‌سازی‌شده از یک فایل پیکربندی کنترلرها برای اعمال سرعت به چرخ‌هااستفاده شده است. این کنترل‌کننده فرمان سرعت را به‌صورت یک بردار چهارتایی دریافت کرده و آن را به مفاصل چرخ‌ها اعمال می‌کند.


\begin{figure}[ht]
	\centering
	\includegraphics[
	width=8cm,
	trim=0 0 0 0,
	clip
	]{Skid-Steer-crop}
	\caption{مقایسه حرکت دیفرانسیلی و حرکت مبتنی بر \lr{skid-steering} در ربات‌های چرخ‌دار \cite{shamah1999effect}}
	\label{fig:skid-steer-illustration}
	
	\vspace{0.4em}
	
	{\small
		\href{https://mohammadabedi1179.github.io/thesis-archive/chapters/ch4/Skid-Steer-crop.html}
		{\textbf{مشاهده نسخه آنلاین شکل}}
		\qquad
		\qrcode[height=1.2cm]{https://mohammadabedi1179.github.io/thesis-archive/chapters/ch4/Skid-Steer-crop.html}
	}
\end{figure}

\begin{figure}[ht]

\end{figure}

در پیکربندی مدل، مفاصل کنترل‌شونده به‌ترتیب زیر تعریف شده‌اند:

\begin{equation}
	\mathbf{q}_{w} =
	\left[
	FL\_joint,\;
	BL\_joint,\;
	FR\_joint,\;
	BR\_joint
	\right]^T
	\label{eq:wheel_joint_order}
\end{equation}

بنابراین، بردار فرمان سرعت ارسالی به کنترل‌کننده نیز باید با همین ترتیب ساخته شود. برای یک ربات چهارچرخ محرک با چرخ‌های محور ثابت، رابطه تقریبی میان سرعت خطی ربات، سرعت زاویه‌ای ربات و سرعت زاویه‌ای چرخ‌های دو سمت به‌صورت زیر در نظر گرفته می‌شود:

\begin{equation}
	v = \frac{r}{2}\left(\omega_r + \omega_l\right),
	\quad
	\omega = \frac{r}{L}\left(\omega_r - \omega_l\right)
	\label{eq:skid_velocities}
\end{equation}

که در آن \(v\) سرعت خطی ربات، \(\omega\) سرعت زاویه‌ای حول محور عمودی، \(r\) شعاع چرخ، \(L\) فاصله مؤثر میان چرخ‌های سمت چپ و راست، و \(\omega_l\) و \(\omega_r\) به‌ترتیب سرعت زاویه‌ای چرخ‌های سمت چپ و راست هستند. در مدل حاضر، شعاع چرخ‌ها برابر با \(r=0.0925~m\) و فاصله عرضی میان چرخ‌های چپ و راست، بر اساس موقعیت هندسی چرخ‌ها نسبت به لینک پایه، برابر با \(L=0.66108~m\) است.

با حل روابط فوق نسبت به سرعت چرخ‌های دو سمت، سرعت زاویه‌ای چرخ‌های راست و چپ به‌ترتیب از روابط زیر به‌دست می‌آید:

\begin{equation}
	\omega_r = \frac{v}{r} + \frac{L}{2r}\omega,
	\quad
	\omega_l = \frac{v}{r} - \frac{L}{2r}\omega
	\label{eq:wheel_velocities}
\end{equation}

در ربات چهارچرخ مورد استفاده، سرعت محاسبه‌شده برای هر سمت به چرخ‌های جلو و عقب همان سمت اعمال می‌شود. بنابراین، نگاشت فرمان سرعت به مفاصل چرخ‌ها به‌صورت زیر تعریف می‌شود:

\begin{equation}
	\omega_{FL} = \omega_{BL} = \omega_l,
	\qquad
	\omega_{FR} = \omega_{BR} = \omega_r
	\label{eq:four_wheel_velocity_mapping}
\end{equation}

با توجه به ترتیب مفاصل در کنترل‌کننده، بردار فرمان نهایی ارسالی به شکل زیر است:

\begin{equation}
	\mathbf{u}_{w} =
	\left[
	\omega_l,\;
	\omega_l,\;
	\omega_r,\;
	\omega_r
	\right]^T
	\label{eq:wheel_velocity_command_vector}
\end{equation}

پیش از ارسال فرمان به کنترل‌کننده، مقادیر سرعت چرخ‌ها در محدوده مجاز محدود می‌شوند تا از تولید فرمان‌های غیرواقعی و ناپایدار در شبیه‌سازی جلوگیری شود. این محدودسازی به‌ویژه در آموزش یادگیری تقویتی اهمیت دارد، زیرا عامل ممکن است در مراحل ابتدایی آموزش فرمان‌های شدید یا نوسانی تولید کند.

روابط ارائه‌شده یک تقریب سینماتیکی برای تبدیل فرمان \((v,\omega)\) به سرعت چرخ‌ها هستند. در حرکت مبتنی بر لغزش، به‌ویژه هنگام چرخش، لغزش جانبی چرخ‌ها نقش مهمی در رفتار واقعی ربات دارد. در مدل شبیه‌سازی‌شده، اثرات تماس، اصطکاک و لغزش از طریق مدل فیزیکی \lr{Gazebo} و پارامترهای تماس چرخ‌ها با زمین لحاظ می‌شوند. بنابراین، لایه کنترل حرکتی نقش واسط میان خروجی سیاست یادگیری تقویتی و مفاصل چرخ‌های ربات را بر عهده دارد و فرمان سطح بالای عامل را به فرمان‌های قابل اعمال در مدل دینامیکی تبدیل می‌کند.

%==========================================================
\section{لایه ادراک و استخراج ویژگی‌ها}
\label{sec:perception_layer}

لایه ادراک وظیفه تبدیل داده‌های خام حسگرها به مجموعه‌ای از ویژگی‌های ساخت‌یافته و کم‌بعد را بر عهده دارد؛ به‌گونه‌ای که این ویژگی‌ها بتوانند به‌عنوان ورودی عامل یادگیری تقویتی مورد استفاده قرار گیرند. از آنجا که مسئله ناوبری مورد بررسی شامل هدف متحرک، موانع دینامیکی و احتمال برخورد در فاصله‌های نزدیک است، سامانه ادراک باید علاوه بر تخمین موقعیت نسبی هدف و موانع، اطلاعاتی درباره فاصله‌های نزدیک‌برد و سرعت نسبی موانع نیز فراهم کند.

در این پژوهش، لایه ادراک از سه زیرسامانه اصلی تشکیل شده است:
\begin{enumerate}
	\item زیرسامانه بینایی مبتنی بر دوربین استریو و آشکارساز \lr{YOLO} برای تشخیص موانع و تخمین موقعیت نسبی آن‌ها؛
	\item گره اختصاصی تشخیص نشانگر هدف برای برآورد فاصله، زاویه و سرعت هدف نسبت به ربات؛
	\item زیرسامانه حسگرهای فراصوت برای دریافت اطلاعات نزدیک‌برد و افزایش ایمنی در فاصله‌های بحرانی.
\end{enumerate}

خروجی این زیرسامانه‌ها پس از پالایش و یکپارچه‌سازی، در قالب ویژگی‌هایی مانند فاصله و زاویه هدف، موقعیت نسبی موانع، سرعت نسبی موانع و فاصله‌های نزدیک‌برد در بردار مشاهده عامل یادگیری تقویتی قرار می‌گیرد.

%----------------------------------------------------------
\subsection{زیرسامانه بینایی استریو و تولید داده‌های عمقی}
\label{subsec:stereo_depth}

برای تخمین فاصله هدف و موانع، دو دوربین جلویی تصاویر هم‌زمان را فراهم می‌کنند و عمق از «اختلاف مکانی%
\LTRfootnote{Disparity}» %
نقاط متناظر به‌دست می‌آید. تصاویر کالیبره‌شده ابتدا از اعوجاج اصلاح و با «تصحیح استریو%
\LTRfootnote{Stereo Rectification}» %
هم‌راستا می‌شوند تا جست‌وجوی تطابق به امتداد افقی هر سطر محدود شود. اختلاف مکانی یک نقطه برابر است با:

\begin{equation}
	d = u_L - u_R
	\label{eq:stereo_disparity_main}
\end{equation}

که در آن \(u_L\) و \(u_R\) مختصات افقی نقطه در تصاویر چپ و راست‌اند. عمق پیکسل‌های معتبر سپس از رابطه زیر محاسبه می‌شود:

\begin{equation}
	Z = \frac{fB}{d}
	\label{eq:stereo_depth_basic}
\end{equation}

در رابطه \ref{eq:stereo_depth_basic}، \(Z\) عمق، \(f\) فاصله کانونی پیکسلی، \(B\) فاصله مبنا و \(d\) اختلاف مکانی است. عمق با اختلاف مکانی رابطه معکوس دارد و مقادیر صفر، منفی یا خارج از بازه معتبر از محاسبات حذف می‌شوند. مطابق شکل \ref{fig:point_cloud_generation}، پس از اصلاح و تطبیق تصاویر، نقشه اختلاف مکانی به عمق تبدیل می‌شود و در صورت نیاز با بازفرافکنی پیکسل‌های معتبر، ابرنقاط تشکیل می‌گردد.
\begin{figure}[ht]
	\centering
	\includegraphics[width=13cm]{point_cloud}
	\caption{فرایند استخراج نقشه عمق و ابرنقاط با استفاده از سامانه
		بینایی استریو}
	\label{fig:point_cloud_generation}
	
	\vspace{0.4em}
	
	{\small
		\href{https://mohammadabedi1179.github.io/thesis-archive/chapters/ch4/point_cloud.html}
		{\textbf{مشاهده نسخه تعاملی شکل با کیفیت بالاتر}}
		\qquad
		\qrcode[height=1.2cm]{https://mohammadabedi1179.github.io/thesis-archive/chapters/ch4/point_cloud.html}
	}
\end{figure}
\begin{figure}[ht]
	
\end{figure}

کیفیت اختلاف مکانی به بافت، روشنایی، سطوح یکنواخت، بازتاب، انسداد و تنظیمات تطبیق وابسته است و نواحی کم‌بافت ممکن است فاقد مقدار معتبر باشند. برای بررسی بهبود خروجی، دو روش به کار رفت: بهبود کنتراست تطبیقی محدودشده%
\LTRfootnote{Contrast Limited Adaptive Histogram Equalization (CLAHE)} %
پیش از تطبیق، که بدون تغییر مستقیم اختلاف مکانی کنتراست موضعی تصاویر را افزایش می‌دهد؛ و فیلتر کمترین مربعات وزن‌دار%
\LTRfootnote{Weighted Least Squares (WLS)} %
پس از تولید نقشه‌های چپ و راست، که نویزهای پراکنده را با حفظ مرزها کاهش می‌دهد. حالت ترکیبی \lr{CLAHE}+\lr{WLS} نیز ارزیابی شد. شکل \ref{fig:disparity_filter_comparison} خروجی خام و سه حالت پردازشی را مقایسه می‌کند.

\begin{figure}[ht]
	\centering
	
	\subfloat[]{
		\includegraphics[width=0.47\textwidth]{Disparity_map_raw.png}
	}
	\hfill
	\subfloat[]{
		\includegraphics[width=0.47\textwidth]{Disparity_map_clahe.png}
	}
	
	\vspace{0.3cm}
	
	\subfloat[]{
		\includegraphics[width=0.47\textwidth]{Disparity_map_wls.png}
	}
	\hfill
	\subfloat[]{
		\includegraphics[width=0.47\textwidth]{Disparity_map_both.png}
	}
	
	\caption{مقایسه کیفی نقشه اختلاف مکانی در حالت‌های مختلف:
		آ) نقشه اختلاف مکانی خام،
		ب) نقشه اختلاف مکانی حاصل از اعمال \lr{CLAHE} بر تصاویر ورودی،
		ج) نقشه اختلاف مکانی پالایش‌شده با فیلتر \lr{WLS}،
		د) نقشه اختلاف مکانی حاصل از ترکیب \lr{CLAHE} و \lr{WLS}.
		نواحی سیاه نشان‌دهنده پیکسل‌هایی هستند که برای آن‌ها مقدار اختلاف
		مکانی معتبر محاسبه نشده است.}
	\label{fig:disparity_filter_comparison}
\end{figure}

همان‌گونه که در شکل \ref{fig:disparity_filter_comparison} مشاهده می‌شود، در نقشه خام، نواحی کم‌بافت، سطوح یکنواخت و مرزهای انسداد پیکسل‌های نامعتبر و پراکنده بیشتری دارند. \lr{CLAHE} جزئیات قابل تطبیق را افزایش می‌دهد، اما ممکن است نویز و تغییرات کوچک شدت را نیز تقویت کند. \lr{WLS} مقادیر پراکنده را کاهش داده و مرز اجسام را منسجم‌تر می‌کند، ولی نمی‌تواند برای تمام نواحی فاقد تطابق، عمق صحیح بسازد. حالت ترکیبی بخشی از نویز را حذف و پیوستگی سطوح را افزایش می‌دهد؛ بااین‌حال، پیوستگی ظاهری الزاماً معادل دقت بیشتر نیست.

رنگ‌های شکل \ref{fig:disparity_filter_comparison} فقط نگاشت اختلاف مکانی‌اند و فاصله متری را نشان نمی‌دهند؛ فاصله با پارامترهای کالیبراسیون و رابطه \ref{eq:stereo_depth_basic} محاسبه می‌شود و نواحی سیاه بیانگر نبود تطابق معتبرند، نه فاصله بسیار زیاد. برای استخراج عمق هر هدف یا مانع، ناحیه متناظر در نقشه انتخاب، پیکسل‌های نامعتبر حذف و یک مقدار نماینده از پیکسل‌های معتبر محاسبه می‌شود. نمونه‌برداری ناحیه‌ای نسبت به یک پیکسل در برابر نویز مقاوم‌تر است، اما حضور پس‌زمینه، چند عمق یا نسبت بالای پیکسل نامعتبر همچنان خطا ایجاد می‌کند؛ ازاین‌رو، اعتبار و نسبت پیکسل‌های معتبر کنترل می‌شوند.

مقایسه ارائه‌شده در شکل \ref{fig:disparity_filter_comparison} ماهیتی کیفی دارد و انتخاب نهایی باید با خطای فاصله، خطای میانگین مربعات، درصد پیکسل معتبر و پایداری در حرکت سنجیده شود؛ بنابراین، دقت فاصله‌سنجی نیز به‌صورت تجربی ارزیابی شده است.

جزئیات هندسه استریو، اصلاح اعوجاج، تصحیح، هندسه اپی‌پولار، تطبیق بلوکی، \lr{CLAHE}، \lr{WLS} و تبدیل اختلاف مکانی به عمق و ابرنقاط در پیوست \ref{appendix:stereo_depth_sgbm} آمده است؛ تمرکز این بخش بر کاربرد خروجی عمقی در تولید ویژگی‌های عامل است.%----------------------------------------------------------
\subsection{آشکارسازی موانع با \lr{YOLO}}
\label{subsec:yolo_obstacle_detection}

برای آشکارسازی موانع دینامیکی در تصویر، از یک گره مستقل مبتنی بر مدل \lr{YOLO} استفاده شده است. این گره تصویر خام دوربین چپ را دریافت کرده و برای هر جسم آشکارشده، یک جعبه محدودکننده دوبعدی به‌همراه کلاس، میزان اطمینان و شناسه جسم منتشر می‌کند. پردازش تصویر به‌صورت مستقیم روی آرایه تصویری انجام می‌شود و از ذخیره‌سازی موقت تصویر روی دیسک استفاده نمی‌گردد؛ این موضوع باعث کاهش تأخیر زنجیره ادراک می‌شود.

در پیکربندی مورد استفاده، گره تشخیص به‌صورت پیش‌فرض تصویر دوربین چپ را پردازش می‌کند؛ زیرا در زنجیره تخمین عمق، جعبه‌های محدودکننده تصویر چپ برای نمونه‌برداری از نقشه اختلاف مکانی کافی هستند. 	در سناریوهای این پژوهش انسان به‌عنوان مانع دینامیکی اصلی در نظر گرفته شده است.

پارامترهای اصلی گره آشکارسازی در جدول \ref{tab:yolo_detector_params} خلاصه شده‌اند.

\begin{table}[ht]
	\centering
	\caption{پارامترهای اصلی گره آشکارسازی مبتنی بر \lr{YOLO}}
	\label{tab:yolo_detector_params}
	\begin{tabular}{|c|c|}
		\hline
		\textbf{پارامتر} & \textbf{مقدار مورد استفاده} \\
		\hline
		مدل آشکارساز & \lr{yolo11} \\
		\hline
		اندازه تصویر ورودی & \(640\) \\
		\hline
		آستانه اطمینان در گره آشکارساز & \(0.25\) \\
		\hline
		حداکثر نرخ پردازش & \(15~Hz\) \\
		\hline
		حالت ردیابی & فعال \\
		\hline
		ردیاب مورد استفاده & \lr{BoT-SORT} \\
		\hline
	\end{tabular}
\end{table}

استفاده از حالت ردیابی باعث می‌شود در صورت امکان، برای هر مانع شناسه‌ای پایدار در فریم‌های متوالی تولید شود. این شناسه در مراحل بعدی برای اتصال اندازه‌گیری‌های زمانی و تخمین سرعت نسبی موانع مورد استفاده قرار می‌گیرد.

%----------------------------------------------------------
\subsection{تخمین عمق و موقعیت دوبعدی موانع از نقشه اختلاف مکانی}
\label{subsec:obstacle_depth_position}

پس از آشکارسازی موانع در تصویر چپ، جعبه محدودکننده هر جسم با نقشه اختلاف مکانی هم‌زمان‌سازی می‌شود. در این مرحله، هدف آن است که از ناحیه دوبعدی آشکارشده در تصویر، یک برآورد عمق و موقعیت نسبی در صفحه حرکت ربات استخراج شود.

خروجی آشکارساز برای هر مانع شامل مرکز جعبه \((u_c,v_c)\)، عرض \(w\) و ارتفاع \(h\) است. بر این اساس، مختصات جعبه محدودکننده در تصویر به‌صورت زیر تعریف می‌شود:

\begin{equation}
	\begin{aligned}
		u_{\min} &= u_c - \frac{w}{2} &
		u_{\max} &= u_c + \frac{w}{2} \\
		v_{\min} &= v_c - \frac{h}{2} &
		v_{\max} &= v_c + \frac{h}{2}
	\end{aligned}
	\label{eq:bbox_limits}
\end{equation}

برای کاهش اثر پیکسل‌های مرزی و نواحی پس‌زمینه، جعبه محدودکننده پیش از نمونه‌برداری از نقشه اختلاف مکانی، از هر طرف کوچک می‌شود. اگر جعبه اولیه به‌صورت \((x_1,y_1,x_2,y_2)\) تعریف شود و ضریب کوچک‌سازی ناحیه موردنظر را \(\alpha\) بنامیم، جعبه کوچک‌شده به‌صورت زیر به‌دست می‌آید:

\begin{equation}
	\begin{aligned}
		x'_1 &= x_1 + \alpha(x_2-x_1) &
		x'_2 &= x_2 - \alpha(x_2-x_1) \\
		y'_1 &= y_1 + \alpha(y_2-y_1) &
		y'_2 &= y_2 - \alpha(y_2-y_1)
	\end{aligned}
	\label{eq:roi_shrink}
\end{equation}

در این پژوهش، این کوچک‌سازی به‌منظور حذف پیکسل‌های ناپایدار اطراف مرز جسم انجام می‌شود؛ زیرا در لبه‌های جعبه محدودکننده معمولاً ترکیبی از پیکسل‌های جسم و پس‌زمینه وجود دارد و این موضوع می‌تواند برآورد اختلاف مکانی را ناپایدار کند.

پس از تعیین ناحیه داخلی، تنها پیکسل‌هایی معتبر در نظر گرفته می‌شوند
که مقدار اختلاف مکانی آن‌ها متناهی و بزرگ‌تر از حداقل اختلاف مکانی
تعیین‌شده باشد. اگر تعداد یا نسبت پیکسل‌های معتبر از حدود مشخص‌شده
کمتر باشد، اندازه‌گیری مربوط به آن جسم در آن فریم کنار گذاشته می‌شود.
مجموعه پیکسل‌های معتبر داخل ناحیه جسم \(i\) با
\(\Omega_i^{valid}\) نشان داده می‌شود. برای کاهش حساسیت تخمین نسبت
به مقادیر پرت، خطاهای تطبیق و پیکسل‌های باقی‌مانده از پس‌زمینه،
میانه مقادیر اختلاف مکانی معتبر به‌عنوان مقدار نماینده ناحیه انتخاب
می‌شود:

\begin{equation}
	\tilde{d}_i =
	\underset{(u,v)\in\Omega_i^{valid}}{\operatorname{median}}
	\left\{d(u,v)\right\}
	\label{eq:median_valid_disparity}
\end{equation}

که در آن \(N_i\) تعداد پیکسل‌های معتبر داخل ناحیه است. سپس عمق خام جسم با استفاده از رابطه استریو به‌صورت زیر محاسبه می‌شود:

\begin{equation}
	Z_i^{\mathrm{raw}} =
	\frac{fB}{\tilde{d}_i}
	\label{eq:obstacle_raw_depth}
\end{equation}

نکته مهم آن است که در این روش، ابتدا مقدار اختلاف مکانی معتبر در ناحیه جسم میانگین‌گیری می‌شود و سپس عمق از روی مقدار میانگین اختلاف مکانی محاسبه می‌گردد. این انتخاب باعث می‌شود تخمین عمق نسبت به نویز پیکسل‌به‌پیکسل در نقشه اختلاف مکانی پایدارتر شود.

برای محاسبه زاویه نسبی جسم نسبت به محور دید دوربین، از مختصات افقی مرکز جعبه محدودکننده استفاده می‌شود:

\begin{equation}
	\theta_i =
	\tan^{-1}
	\left(
	\frac{u_{c,i}-c_x}{f_x}
	\right)
	\label{eq:obstacle_bearing}
\end{equation}

که در آن \(u_{c,i}\) مختصات افقی مرکز جعبه، \(c_x\) مختصات مرکز اپتیکی و \(f_x\) فاصله کانونی افقی دوربین است. با فرض آن‌که \(Z_i\) فاصله رو به جلو در چارچوب دوربین را نشان دهد، مختصات تقریبی جسم در صفحه حرکت به‌صورت زیر به‌دست می‌آید:

\begin{equation}
	x_{\text{cam},i} = Z_i,
	\qquad
	y_{\text{cam},i} = Z_i \tan(\theta_i)
	\label{eq:camera_xy_from_depth}
\end{equation}

سپس با لحاظ کردن مکان نصب دوربین نسبت به لینک پایه، مختصات جسم در چارچوب ربات محاسبه می‌شود. در پیاده‌سازی حاضر، مؤلفه طولی دوربین نسبت به لینک پایه به مختصات \(x\) افزوده شده و مؤلفه جانبی با توجه به قرارداد محورها اصلاح می‌گردد:

\begin{equation}
	x_i = x_{\text{cam},i} + x_{\text{cam-offset}},
	\qquad
	y_i = y_{\text{cam-offset}} - y_{\text{cam},i}
	\label{eq:base_link_obstacle_xy}
\end{equation}

در این رابطه، \(x_i\) و \(y_i\) مختصات نسبی مانع در چارچوب لینک پایه هستند. این مختصات مبنای محاسبه فاصله، زاویه نسبی و سرعت مانع در مراحل بعدی قرار می‌گیرند.

\subsubsection{ارزیابی تجربی دقت تخمین عمق و موقعیت موانع در شبیه سازی}

برای ارزیابی عملکرد سامانه، عمق استخراج‌شده از نقشه اختلاف مکانی با فاصله واقعی موانع نسبت به دوربین مقایسه شد. در این ارزیابی، نتایج مربوط به فریم‌هایی هستند که آشکارسازی مانع و استخراج اختلاف مکانی معتبر با موفقیت انجام شده است.

شکل \ref{fig:real_depth_estimation_accuracy} مقایسه عمق واقعی و عمق تخمین‌زده‌شده را در حالت خام و پالایش‌شده نشان می‌دهد. نزدیکی نقاط به خط ایده آل (y=x) بیانگر دقت بیشتر تخمین است.

\begin{figure}[htb]
	\centering
	\includegraphics[width=0.75\linewidth]{01_depth_scatter_obstacle.png}
	\caption{مقایسه عمق واقعی موانع با عمق تخمین‌زده‌شده در حالت خام و
		پالایش‌شده؛ خط‌چین نشان‌دهنده حالت ایده آل
		\(Z_{\mathrm{estimated}}=Z_{\mathrm{true}}\) است}
	\label{fig:real_depth_estimation_accuracy}
\end{figure}

نتایج نشان می‌دهند که تخمین عمق خام، روند تغییر فاصله واقعی را با دقت مناسبی دنبال می‌کند و ضریب تعیین آن برابر با ($R^2=0.991$) است. پالایش زمانی اگرچه بخشی از نوسانات و بایاس میانگین را کاهش می‌دهد، به‌دلیل ایجاد تأخیر در پاسخ، الزاماً موجب بهبود دقت مطلق نمی‌شود. این تأخیر به‌ویژه هنگامی اهمیت پیدا می‌کند که فاصله نسبی مانع در فریم‌های متوالی تغییر کند.

برای بررسی اثر حرکت اسکوتر، تغییر خطای تخمین عمق نسبت به سرعت طولی و سرعت زاویه‌ای در شکل \ref{fig:depth_error_motion_effects} ارائه شده است.

\begin{figure}[htbp]
	\centering
	
	\subfloat[]{
		\includegraphics[width=0.48\textwidth]
		{02_error_vs_forward_speed_obstacle.png}
	}
	\hfill
	\subfloat[]{
		\includegraphics[width=0.48\textwidth]
		{03_error_vs_angular_speed_obstacle.png}
	}
	
	\caption{بررسی اثر حرکت اسکوتر بر خطای تخمین عمق:
		الف) تغییر خطای علامت‌دار و قدر مطلق خطا نسبت به سرعت طولی،
		ب) تغییر خطای علامت‌دار و قدر مطلق خطا نسبت به سرعت زاویه‌ای}
	\label{fig:depth_error_motion_effects}
	
\end{figure}

به‌طور کلی، افزایش سرعت حرکت می‌تواند به کاهش پایداری تخمین منجر شود. این اثر در سرعت زاویه‌ای اهمیت بیشتری دارد؛ زیرا دوران اسکوتر موجب جابه‌جایی سریع مانع در تصویر، تغییر ناحیه مورد بررسی و کاهش کیفیت تطبیق نقاط متناظر می‌شود. بااین‌حال، به‌دلیل ایدئال‌بودن تصاویر و هماهنگی مناسب حسگرها در محیط شبیه‌سازی، افزایش خطا با سرعت در نتایج حاضر چندان شدید و یکنواخت نیست.

علاوه بر حرکت ربات، فاصله مانع و موقعیت آن در میدان دید نیز بر دقت تخمین اثر می‌گذارند. شکل \ref{fig:depth_error_depth_bearing} تغییرات خطای عمق را نسبت به فاصله واقعی و زاویه قرارگیری مانع نشان می‌دهد.

\begin{figure}[htb]
	\centering
	\includegraphics[width=0.75\linewidth]
	{04_error_vs_depth_and_bearing_obstacle.png}
	\caption{تغییر خطای علامت‌دار تخمین عمق بر حسب:
		شکل سمت راست) فاصله واقعی مانع از دوربین،
		شکل سمت چپ) زاویه واقعی مانع نسبت به محور اپتیکی دوربین}
	\label{fig:depth_error_depth_bearing}
\end{figure}

با افزایش فاصله مانع، مقدار اختلاف مکانی کوچک‌تر شده و حساسیت محاسبه عمق نسبت به خطاهای پیکسلی افزایش می‌یابد. این رفتار از رابطه استریو قابل توضیح است:

\begin{equation}
	\Delta Z
	\approx
	\frac{Z^2}{fB}\Delta d
	\label{eq:stereo_depth_error_propagation}
\end{equation}

بر این اساس، یک خطای تقریباً ثابت در اختلاف مکانی، در فواصل بیشتر به خطای بزرگ‌تری در تخمین عمق تبدیل می‌شود. همچنین، در زاویه‌های بزرگ‌تر و نزدیک لبه‌های تصویر، کاهش هم‌پوشانی دید دو دوربین، اعوجاج باقی‌مانده و کوچک‌شدن ناحیه معتبر جسم می‌تواند موجب افزایش پراکندگی خطا شود.

در مجموع، بهترین عملکرد زمانی حاصل می‌شود که مانع در فاصله نزدیک‌تر، ناحیه مرکزی میدان دید و شرایط حرکت آرام قرار داشته باشد. با افزایش فاصله و سرعت ربات، به‌ویژه سرعت زاویه‌ای، دقت تخمین کاهش می‌یابد؛ هرچند به‌دلیل شرایط ایده آل محیط شبیه‌سازی، شدت این کاهش در نتایج حاضر چندان محسوس نیست.


%----------------------------------------------------------
\subsection{پالایش زمانی و تخمین سرعت موانع با فیلتر کالمن}
\label{subsec:obstacle_kalman_velocity}

نویز تطبیق، پیکسل‌های نامعتبر، مرز جعبه و تغییرات آشکارساز می‌توانند خروجی ادراک و فرمان عامل را ناپایدار کنند. ازاین‌رو، پالایش دومرحله‌ای شامل حذف نقاط پرت با «بازه میان‌چارکی%
\LTRfootnote{Interquartile Range-IQR}» %
و تخمین موقعیت و سرعت نسبی با کالمن استفاده شد. برای هر مانع، تاریخچه‌ای مستقل با شناسه \lr{YOLO} یا کلید جایگزین مبتنی بر کلاس و مرکز گسسته‌شده جعبه نگه‌داری می‌شود. اختلاف‌های اخیر در بافر قرار می‌گیرند و \lr{IQR} به‌صورت زیر محاسبه می‌شود:

\begin{equation}
	\mathrm{IQR}=Q_3-Q_1
	\label{eq:iqr_definition}
\end{equation}

نمونه‌های خارج از بازه زیر پرت در نظر گرفته می‌شوند:

\begin{equation}
	d_{\mathrm{low}} = Q_1 - 1.5\,\mathrm{IQR},
	\qquad
	d_{\mathrm{high}} = Q_3 + 1.5\,\mathrm{IQR}
	\label{eq:iqr_bounds}
\end{equation}

پس از حذف نقاط پرت، میانگین داده‌های باقی‌مانده استفاده می‌شود؛ در صورت ناکافی‌بودن نمونه‌ها، مقدار خام جاری حفظ می‌گردد. سپس مختصات دوبعدی به فیلتر کالمن مستقل با مدل سرعت ثابت و بردار حالت زیر داده می‌شود:

\begin{equation}
	\mathbf{x}_k =
	\begin{bmatrix}
		x_k \\
		v_{x,k} \\
		y_k \\
		v_{y,k}
	\end{bmatrix}
	\label{eq:kalman_state_obstacle}
\end{equation}

\(x_k,y_k\) موقعیت و \(v_{x,k},v_{y,k}\) سرعت نسبی در چارچوب پایه‌اند. فقط موقعیت اندازه‌گیری می‌شود و سرعت از تغییرات زمانی آن با مدل گذار زیر تخمین زده می‌شود:

\begin{equation}
	\mathbf{F}_k =
	\begin{bmatrix}
		1 & \Delta t & 0 & 0 \\
		0 & 1        & 0 & 0 \\
		0 & 0        & 1 & \Delta t \\
		0 & 0        & 0 & 1
	\end{bmatrix}
	\label{eq:kalman_transition_obstacle}
\end{equation}

\(\Delta t\) از زمان واقعی دو به‌روزرسانی و در بازه‌ای محدود محاسبه می‌شود. ماتریس‌های نویز، مقداردهی اولیه و مراحل فیلتر در پیوست \ref{appendix:kalman_speed} آمده‌اند.

خروجی فیلتر، موقعیت پالایش‌شده \((x_i,y_i)\) و سرعت نسبی \((v_{x,i},v_{y,i})\) است. اندازه سرعت از رابطه زیر به‌دست می‌آید:

\begin{equation}
	v_i =
	\sqrt{
		v_{x,i}^{2}
		+
		v_{y,i}^{2}
	}
	\label{eq:obstacle_speed}
\end{equation}
این سرعت در چارچوب پایه نسبی است و سرعت مطلق جهانی مانع را نشان نمی‌دهد.
%--------------------------------------------------------
\subsection{تشخیص هدف و برآورد فاصله و زاویه نشانگر}
\label{subsec:goal_marker_detection}

در سامانه ناوبری حاضر، هدف توسط یک نشانگر دیداری مشخص در محیط شبیه‌سازی تعریف شده است. این نشانگر دارای نواحی رنگی قابل تشخیص، به‌ویژه بخش‌های قرمز، است تا بتوان آن را در تصویر دوربین چپ به‌صورت پایدار شناسایی کرد. برای استخراج موقعیت نسبی هدف، یک گره مستقل طراحی شده است که تصویر دوربین چپ و نقشه اختلاف مکانی را به‌صورت هم‌زمان دریافت می‌کند.

\subsubsection{تشخیص رنگی هدف}

تشخیص هدف بر پایه استخراج نواحی قرمز تصویر انجام می‌شود. ابتدا تصویر دوربین چپ از فضای رنگی \lr{BGR} به فضای \lr{HSV} تبدیل می‌شود. سپس دو بازه رنگی متناظر با رنگ قرمز برای ساخت ماسک دودویی استفاده می‌گردد. پس از تشکیل ماسک، عملیات مورفولوژیکی بازشدن و بسته‌شدن برای کاهش نویز و پرکردن نواحی گسسته اعمال می‌شود. در نهایت، بزرگ‌ترین کانتور معتبر به‌عنوان جعبه محدودکننده هدف انتخاب می‌گردد.

برای افزایش سرعت و پایداری تشخیص، اگر هدف در فریم قبلی مشاهده شده باشد، ابتدا جست‌وجو در اطراف آخرین جعبه محدودکننده انجام می‌شود. در صورت شکست جست‌وجوی محلی، کل تصویر برای یافتن هدف بررسی می‌گردد. اگر هیچ ناحیه معتبری یافت نشود، هدف نامرئی فرض شده و خروجی گره با مقدار عدم شناسایی منتشر می‌شود.

\subsubsection{محاسبه عمق هدف از نقشه اختلاف مکانی}

پس از تشخیص جعبه هدف، ناحیه داخلی جعبه به‌منظور کاهش اثر لبه‌ها کوچک می‌شود. سپس پیکسل‌های معتبر اختلاف مکانی داخل این ناحیه استخراج می‌شوند. برخلاف مسیر پردازش موانع که برای هر جسم میانگین اختلاف مکانی محاسبه می‌شود، در گره هدف از میانه اختلاف مکانی استفاده شده است:

\begin{equation}
	d_{\mathrm{goal}} =
	\mathrm{median}
	\left(
	d(u,v) \; ; \; (u,v)\in\Omega_{\mathrm{goal}}^{valid}
	\right)
	\label{eq:goal_median_disparity}
\end{equation}

استفاده از میانه باعث کاهش اثر پیکسل‌های پرت، خطاهای مرزی جعبه و جهش‌های موضعی نقشه اختلاف مکانی می‌شود. سپس فاصله خام هدف نسبت به دوربین از رابطه زیر به‌دست می‌آید:

\begin{equation}
	Z_{\mathrm{goal}} =
	\frac{fB}{d_{\mathrm{goal}}}
	\label{eq:goal_raw_depth}
\end{equation}

برای جلوگیری از پذیرش مقادیر غیرواقعی، تنها اندازه‌گیری‌هایی پذیرفته می‌شوند که مقدار اختلاف مکانی آن‌ها از حداقل معتبر بزرگ‌تر بوده و عمق محاسبه‌شده در بازه فیزیکی قابل قبول قرار داشته باشد.

\subsubsection{محاسبه زاویه و تبدیل به چارچوب لینک پایه}

زاویه نسبی هدف نسبت به محور دید دوربین از روی مختصات افقی مرکز جعبه هدف محاسبه می‌شود:

\begin{equation}
	\theta_{\mathrm{goal}} =
	\tan^{-1}
	\left(
	\frac{u_{\mathrm{goal}}-C_x}{f_x}
	\right)
	\label{eq:goal_camera_bearing}
\end{equation}

که در آن \(u_{\mathrm{goal}}\) مختصات افقی مرکز جعبه هدف، \(C_x\) مرکز اپتیکی تصویر چپ و \(f_x\) فاصله کانونی افقی دوربین است. با توجه به فاصله نصب دوربین نسبت به لینک پایه، مختصات هدف در صفحه حرکت ربات به‌صورت زیر محاسبه می‌شود:

\begin{equation}
	x_{\mathrm{goal}} =
	Z_{\mathrm{goal}} + x_{\mathrm{cam-offset}},
	\qquad
	y_{\mathrm{goal}} =
	Z_{\mathrm{goal}}\tan(\theta_{\mathrm{goal}}) + y_{\mathrm{cam-offset}}
	\label{eq:goal_base_xy}
\end{equation}

در نتیجه، فاصله و زاویه هدف نسبت به چارچوب پایه ربات برابر خواهد بود با:

\begin{equation}
	\rho_{\mathrm{goal}} =
	\sqrt{x_{\mathrm{goal}}^2 + y_{\mathrm{goal}}^2},
	\qquad
	\phi_{\mathrm{goal}} =
	\tan^{-1}
	\left(
	\frac{y_{\mathrm{goal}}}{x_{\mathrm{goal}}}
	\right)
	\label{eq:goal_range_bearing_base}
\end{equation}

\subsubsection{هموارسازی فاصله هدف}

برای کاهش نوسانات لحظه‌ای فاصله هدف، مقدار فاصله اصلاح‌شده با استفاده از فیلتر «میانگین متحرک نمایی%
\LTRfootnote{Exponential Moving Average-EMA}» %
هموار می‌شود. رابطه به‌روزرسانی این فیلتر به‌صورت زیر است:

\begin{equation}
	\rho_{\mathrm{EMA}}^{(k)}
	=
	\alpha_{\mathrm{EMA}}\rho_{\mathrm{goal}}^{(k)}
	+
	\left(1-\alpha_{\mathrm{EMA}}\right)
	\rho_{\mathrm{EMA}}^{(k-1)}
	\label{eq:goal_ema}
\end{equation}

که در آن \(\alpha_{\mathrm{EMA}}\) ضریب هموارسازی است. اگر برای هدف هنوز مقدار قبلی وجود نداشته باشد، مقدار اولیه فیلتر برابر با اولین اندازه‌گیری معتبر قرار داده می‌شود. خروجی نهایی گره هدف شامل وضعیت رؤیت‌پذیری، فاصله هموارشده، زاویه نسبی و مختصات پیکسلی مرکز هدف است و در قالب پیام جداگانه منتشر می‌شود.


\subsubsection{ارزیابی عملکرد در محیط شبیه‌سازی}
\label{subsubsec:sim_goal_marker_evaluation}

برای ارزیابی عملکرد سامانه، فاصله محاسبه‌شده هدف با مقدار مرجع استخراج‌شده از شبیه‌ساز مقایسه شد. شکل \ref{fig:sim_goal_depth_accuracy} مقایسه عمق مرجع و عمق تخمین‌زده‌شده را نشان می‌دهد.

\begin{figure}[htb]
	\centering
	\includegraphics[width=0.5\linewidth]
	{sim_goal_depth_scatter.png}
	\caption{مقایسه عمق مرجع هدف با عمق محاسبه‌شده از نقشه اختلاف
		مکانی در محیط شبیه‌سازی}
	\label{fig:sim_goal_depth_accuracy}
\end{figure}

نتایج نشان می‌دهند که فاصله تخمین‌زده‌شده، روند کلی تغییر فاصله واقعی را به‌خوبی دنبال می‌کند و ضریب تعیین آن برابر با ($R^2=0.991$) است. بااین‌حال، وجود بایاس مثبت نشان می‌دهد که سامانه فاصله هدف را اندکی بیشتر از مقدار واقعی تخمین می‌زند. مقایسه خروجی خام و هموارشده نیز نشان داد که فیلتر \lr{EMA} بیشتر موجب کاهش نوسانات لحظه‌ای شده و بهبود محسوسی در دقت کلی ایجاد نکرده است.

اثر سرعت طولی و زاویه‌ای ربات بر خطای تخمین فاصله در شکل \ref{fig:sim_goal_motion_effects} ارائه شده است.

\begin{figure}[htbp]
	\centering
	\subfloat[]{
		\includegraphics[width=0.48\textwidth]
		{sim_goal_error_vs_forward_speed.png}
	}
	\hfill
	\subfloat[]{
		\includegraphics[width=0.48\textwidth]
		{sim_goal_error_vs_angular_speed.png}
	}
	
	\caption{اثر حرکت ربات بر خطای تخمین فاصله هدف:
		الف) سرعت طولی، ب) سرعت زاویه‌ای}
	\label{fig:sim_goal_motion_effects}
	
\end{figure}

با افزایش سرعت حرکت ربات، پایداری تخمین تا حدی کاهش می‌یابد. این اثر در سرعت زاویه‌ای مهم‌تر است؛ زیرا دوران اسکوتر موجب جابه‌جایی سریع‌تر هدف در تصویر و کاهش ثبات ناحیه تشخیص و نقاط متناظر استریو می‌شود. بااین‌حال، به‌دلیل ایدئال‌بودن تصاویر، نبود نویز شدید و هماهنگی مناسب حسگرها در شبیه‌ساز، این افزایش خطا در نتایج حاضر چندان محسوس نیست.

تغییرات خطا نسبت به فاصله مرجع در شکل \ref{fig:sim_goal_error_vs_depth} ارائه شده است.

\begin{figure}[htbp]
	\centering
	\includegraphics[width=0.82\linewidth]
	{sim_goal_error_vs_true_depth.png}
	\caption{تغییر خطای تخمین فاصله هدف بر حسب عمق مرجع شبیه‌ساز}
	\label{fig:sim_goal_error_vs_depth}
\end{figure}

با افزایش فاصله هدف، اختلاف مکانی کوچک‌تر شده و حساسیت تخمین عمق نسبت به خطاهای پیکسلی افزایش می‌یابد. این موضوع را می‌توان به‌صورت زیر بیان کرد:

\begin{equation}
	\Delta Z
	\approx
	\frac{Z^2}{fB}\Delta d
	\label{eq:sim_goal_depth_uncertainty}
\end{equation}

بنابراین، در فواصل بیشتر، یک خطای محدود در اختلاف مکانی می‌تواند به خطای بزرگ‌تری در فاصله تخمین‌زده‌شده تبدیل شود.

ارزیابی زاویه هدف نیز در شکل \ref{fig:sim_goal_bearing_accuracy} ارائه شده است.

\begin{figure}[htbp]
	\centering
	\includegraphics[width=0.95\linewidth]
	{sim_goal_bearing_analysis.png}
	\caption{مقایسه زاویه محاسبه‌شده هدف (شکل سمت راست)با زاویه مرجع و توزیع خطای
		زاویه در محیط شبیه‌سازی (شکل سمت چپ)}
	\label{fig:sim_goal_bearing_accuracy}
\end{figure}

در زاویه‌های بزرگ‌تر و نزدیک لبه‌های تصویر، پراکندگی خطا افزایش می‌یابد. این رفتار می‌تواند ناشی از کاهش دقت تعیین مرکز جعبه محدودکننده، اعوجاج نواحی کناری و کاهش کیفیت تطبیق استریو باشد.

در مجموع، سامانه در محیط شبیه‌سازی توانسته است فاصله و زاویه هدف را با دقت مناسبی دنبال کند. بیشترین دقت زمانی حاصل می‌شود که هدف در فاصله نزدیک‌تر، بخش مرکزی میدان دید و شرایط حرکت آرام قرار داشته باشد. با افزایش فاصله و سرعت ربات، به‌ویژه سرعت زاویه‌ای، دقت تخمین کاهش می‌یابد؛ اما به‌دلیل شرایط ایده آل محیط شبیه‌سازی، شدت این کاهش در نتایج فعلی محدود است.
%----------------------------------------------------------
\subsection{حسگرهای فراصوت و ایمنی نزدیک‌برد}
\label{subsec:ultrasonic_safety}

اگرچه سامانه بینایی استریو امکان تخمین موقعیت موانع را فراهم می‌کند، اما در فاصله‌های بسیار نزدیک ممکن است به دلیل نواحی کور، خطای اختلاف مکانی یا محدودیت میدان دید، اطلاعات دیداری به‌تنهایی برای تصمیم‌گیری ایمن کافی نباشد. به همین دلیل، در مدل ربات پنج حسگر فراصوت نزدیک‌برد استفاده شده است. این حسگرها شامل دو حسگر جلویی، یک حسگر سمت چپ، یک حسگر سمت راست و یک حسگر عقبی هستند.

حسگرهای فراصوت فاصله نزدیک‌ترین جسم در راستای اندازه‌گیری خود را تولید می‌کنند و به‌عنوان اطلاعات مکمل در بردار مشاهده عامل مورد استفاده قرار می‌گیرند. نقش اصلی آن‌ها افزایش حساسیت سامانه نسبت به موانع نزدیک، کاهش احتمال برخورد در فاصله‌های بحرانی و فراهم‌کردن اطلاعات نزدیک‌برد مستقل از سامانه بینایی است.

%----------------------------------------------------------
\subsection{خروجی نهایی لایه ادراک برای عامل یادگیری تقویتی}
\label{subsec:perception_output_rl}

خروجی لایه ادراک، مجموعه‌ای از ویژگی‌های مربوط به هدف، موانع و فاصله‌های نزدیک‌برد است که پس از سازمان‌دهی در قالب بردار مشاهده، به عامل یادگیری تقویتی داده می‌شود. این ویژگی‌ها شامل موارد زیر هستند:

\begin{itemize}
	\item فاصله و زاویه نسبی هدف نسبت به ربات؛
	\item وضعیت رؤیت‌پذیری هدف؛
	\item موقعیت نسبی موانع در چارچوب لینک پایه؛
	\item فاصله و زاویه نسبی موانع؛
	\item مؤلفه‌های سرعت نسبی موانع؛
	\item داده‌های پاک‌سازی‌شده حسگرهای فراصوت در جهات جلو، چپ، راست و عقب.
\end{itemize}

این ساختار باعث می‌شود عامل یادگیری تقویتی به‌جای کار مستقیم با داده‌های خام تصویری یا حسگری، با یک نمایش فشرده و معنادار از محیط تعامل داشته باشد. چنین نمایشی از یک سو هزینه محاسباتی آموزش را کاهش می‌دهد و از سوی دیگر، اطلاعات لازم برای تصمیم‌گیری ایمن در محیط‌های دینامیکی را در اختیار عامل قرار می‌دهد.
%==========================================================
\section{طراحی عامل یادگیری تقویتی}
\label{sec:rl_agent_design}

در این بخش، فرمول‌بندی نهایی عامل یادگیری تقویتی و اجزای اصلی محیط آموزشی مطابق پیاده‌سازی انجام‌شده تشریح می‌شود. برخلاف بخش‌های نظری فصل قبل که الگوریتم \lr{SAC} را به‌صورت عمومی معرفی کردند، تمرکز این بخش بر نحوه تبدیل داده‌های ادراکی، وضعیت ربات، اطلاعات هدف، موانع و منطق ایمنی به ساختاری است که مستقیماً توسط عامل استفاده می‌شود. در پیاده‌سازی حاضر، مشاهده عامل به‌صورت یک دیکشنری شامل دو بخش اصلی تعریف شده است: بردار ویژگی‌های ثابت و ماتریس ویژگی‌های موانع. سپس این دو بخش توسط سیاست سفارشی \lr{LSTM-SAC} پردازش شده و عمل پیوسته ربات در قالب سرعت خطی و سرعت زاویه‌ای تولید می‌شود.

از آنجا که مسئله ناوبری مورد بررسی دارای عدم‌قطعیت ادراکی، موانع متحرک و هدف مرحله‌ای/متحرک است، در کنار آموزش معمول \lr{SAC}، دو سازوکار تکمیلی نیز در طراحی نهایی لحاظ شده‌اند. نخست، یک «سپر ایمنی غیرفعال%
\LTRfootnote{Passive Safety Shield}» %
 که سطح خطر را تشخیص داده، آن را به مشاهده عامل اضافه کرده و از طریق هزینه لاگرانژی عامل را از نواحی خطر دور می‌کند. دوم، «سازوکار معلم--دانشجو%
 \LTRfootnote{Teacher--Student Distillation}» %
  که در آموزش دانشجو، عمل تولیدشده توسط یک معلم آموزش‌دیده روی هندسه دقیق را به‌عنوان لنگر رفتاری در تابع زیان وارد می‌کند.

%----------------------------------------------------------
%----------------------------------------------------------
\subsection{فضای مشاهده عامل}
\label{subsec:observation_space_final}

عامل به حالت کامل دسترسی ندارد و مشاهده از ادراک، اودومتری، فراصوت، سپر ایمنی و حافظه موانع ساخته می‌شود. برای تفکیک ماهیت داده‌ها، مشاهده به‌صورت دیکشنری زیر تعریف شده است:

\begin{equation}
	\mathcal{O}_t =
	\left\{
	\mathbf{s}^{fixed}_t,\;
	\mathbf{O}_t
	\right\}
	\label{eq:dict_observation_final}
\end{equation}

بردار ثابت مستقیماً و ماتریس موانع پس از پردازش \lr{LSTM} وارد استخراج‌گر ویژگی می‌شوند.

%----------------------------------------------------------
\subsubsection{بردار ویژگی‌های ثابت}
\label{subsubsec:fixed_features_final}

بخش ثابت مشاهده، برداری ۱۵بعدی از اطلاعات هدف، حرکت ربات، پنج حسگر فراصوت، فرمان قبلی سیاست و وضعیت سپر است:

\begin{equation}
	\mathbf{s}^{fixed}_t =
	\begin{bmatrix}
		\mathbf{s}^{goal}_t \\
		\mathbf{s}^{robot}_t \\
		\mathbf{s}^{ultra}_t \\
		\mathbf{s}^{policy}_t \\
		\mathbf{s}^{shield}_t
	\end{bmatrix}
	\in \mathbb{R}^{15}
	\label{eq:fixed_obs_15d_final}
\end{equation}

زیر‌بردار هدف، فاصله، زاویه و دو مؤلفه سرعت نسبی را شامل می‌شود:

\begin{equation}
	\mathbf{s}^{goal}_t =
	\begin{bmatrix}
		\tilde{d}_g &
		\tilde{\phi}_g &
		\tilde{v}^{rel}_{g,x} &
		\tilde{v}^{rel}_{g,y}
	\end{bmatrix}^{T}
	\in \mathbb{R}^{4}
	\label{eq:goal_fixed_features}
\end{equation}

زیر‌بردار حرکت ربات شامل سرعت خطی و زاویه‌ای است:

\begin{equation}
	\mathbf{s}^{robot}_t =
	\begin{bmatrix}
		\tilde{v} &
		\tilde{\omega}
	\end{bmatrix}^{T}
	\in \mathbb{R}^{2}
	\label{eq:robot_fixed_features}
\end{equation}

پنج فاصله فراصوت در زیر‌بردار زیر قرار می‌گیرند:

\begin{equation}
	\mathbf{s}^{ultra}_t =
	\begin{bmatrix}
		\tilde{u}_{FL} &
		\tilde{u}_{FR} &
		\tilde{u}_{LS} &
		\tilde{u}_{RS} &
		\tilde{u}_{B}
	\end{bmatrix}^{T}
	\in \mathbb{R}^{5}
	\label{eq:ultrasonic_fixed_features}
\end{equation}

فرمان خطی و زاویه‌ای سیاست در گام قبل نیز افزوده می‌شود:

\begin{equation}
	\mathbf{s}^{policy}_t =
	\begin{bmatrix}
		\tilde{v}^{policy}_{t-1} &
		\tilde{\omega}^{policy}_{t-1}
	\end{bmatrix}^{T}
	\in \mathbb{R}^{2}
	\label{eq:policy_fixed_features}
\end{equation}

زیر‌بردار سپر، شدت خطر و مدت حضور پیوسته در وضعیت غیرایمن را نگه می‌دارد:

\begin{equation}
	\mathbf{s}^{shield}_t =
	\begin{bmatrix}
		\ell_{shield} &
		\tau_{shield}
	\end{bmatrix}^{T}
	\in \mathbb{R}^{2}
	\label{eq:shield_fixed_features}
\end{equation}

علامت \(\tilde{\cdot}\) مقدار نرمال‌شده را نشان می‌دهد. جدول \ref{tab:fixed_obs_features_final} اجزا و نرمال‌سازی را خلاصه می‌کند.

\begin{table}[ht]
	\centering
	\caption{ساختار بردار ویژگی‌های ثابت مشاهده عامل}
	\label{tab:fixed_obs_features_final}
	\resizebox{\textwidth}{!}{%
		\begin{tabular}{|c|c|c|c|}
			\hline
			\textbf{گروه} & \textbf{ویژگی} & \textbf{معنا} & \textbf{نرمال‌سازی/بازه} \\
			\hline
			\multirow{4}{*}{هدف}
			& \(\tilde{d}_g\) & فاصله هدف تا ربات & \(d_g / 15\) در بازه \([0,1]\) \\
			\cline{2-4}
			& \(\tilde{\phi}_g\) & زاویه نسبی هدف & \(\phi_g / \pi\) در بازه \([-1,1]\) \\
			\cline{2-4}
			& \(\tilde{v}^{rel}_{g,x}\) & سرعت نسبی هدف در راستای طولی & \(v^{rel}_{g,x}/2\) در بازه \([-1,1]\) \\
			\cline{2-4}
			& \(\tilde{v}^{rel}_{g,y}\) & سرعت نسبی هدف در راستای جانبی & \(v^{rel}_{g,y}/2\) در بازه \([-1,1]\) \\
			\hline
			
			\multirow{2}{*}{ربات}
			& \(\tilde{v}\) & سرعت خطی ربات & \(v/1\) در بازه \([-1,1]\) \\
			\cline{2-4}
			& \(\tilde{\omega}\) & سرعت زاویه‌ای ربات & \(\omega/1\) در بازه \([-1,1]\) \\
			\hline
			
			\multirow{5}{*}{فراصوت}
			& \(\tilde{u}_{FL}\) & فاصله فراصوت جلو-چپ & \(u_{FL}/3\) در بازه \([0,1]\) \\
			\cline{2-4}
			& \(\tilde{u}_{FR}\) & فاصله فراصوت جلو-راست & \(u_{FR}/3\) در بازه \([0,1]\) \\
			\cline{2-4}
			& \(\tilde{u}_{LS}\) & فاصله فراصوت سمت چپ & \(u_{LS}/3\) در بازه \([0,1]\) \\
			\cline{2-4}
			& \(\tilde{u}_{RS}\) & فاصله فراصوت سمت راست & \(u_{RS}/3\) در بازه \([0,1]\) \\
			\cline{2-4}
			& \(\tilde{u}_{B}\) & فاصله فراصوت عقب & \(u_B/3\) در بازه \([0,1]\) \\
			\hline
			
			\multirow{2}{*}{فرمان قبلی}
			& \(\tilde{v}^{policy}_{t-1}\) & فرمان خطی سیاست در گام قبل & \(v^{policy}_{t-1}/1\) در بازه \([-1,1]\) \\
			\cline{2-4}
			& \(\tilde{\omega}^{policy}_{t-1}\) & فرمان زاویه‌ای سیاست در گام قبل & \(\omega^{policy}_{t-1}/1\) در بازه \([-1,1]\) \\
			\hline
			
			\multirow{2}{*}{سپر ایمنی}
			& \(\ell_{shield}\) & شدت سپر ایمنی & \(0, 0.33, 0.67, 1\) برای سطوح ایمن، احتیاط، خطرناک و بحرانی \\
			\cline{2-4}
			& \(\tau_{shield}\) & مدت حضور پیوسته در وضعیت سپر & \(\min(N_{shield}/50,1)\) \\
			\hline
		\end{tabular}%
	}
\end{table}

در آموزش معلم، هندسه هدف و موانع دقیق است؛ در آموزش دانشجو، موانع از مسیر تشخیص و استریو وارد می‌شوند. در مسیر غیرمعلم، فاصله و زاویه هدف فعلاً از اودومتری نشانگر محاسبه و پیام تشخیص برای زنجیره ادراک و رؤیت‌پذیری استفاده می‌شود. هنگام نبود داده معتبر، آخرین مشاهده معتبر حفظ می‌گردد.

%----------------------------------------------------------
\subsubsection{ماتریس ویژگی‌های موانع}
\label{subsubsec:obstacle_features_final}

بخش دوم مشاهده، ماتریس موانع با حداکثر قابل تنظیم و مقدار پیش‌فرض ۱۰ است:

\begin{equation}
	\mathbf{O}_t =
	\begin{bmatrix}
		\mathbf{o}_{1,t} \\
		\mathbf{o}_{2,t} \\
		\vdots \\
		\mathbf{o}_{10,t}
	\end{bmatrix}
	\in \mathbb{R}^{10\times 5}
	\label{eq:obstacle_matrix_10x5_final}
\end{equation}

هر ردیف، مانع را با پنج ویژگی نمایش می‌دهد:

\begin{equation}
	\mathbf{o}_{i,t}=
	\begin{bmatrix}
		\tilde{x}_{i,t} &
		\tilde{y}_{i,t} &
		\tilde{d}_{i,t} &
		\tilde{v}^{rel}_{x,i,t} &
		\tilde{v}^{rel}_{y,i,t}
	\end{bmatrix}^{T}
	\label{eq:single_obstacle_feature_final}
\end{equation}

این ویژگی‌ها مختصات پایه، فاصله و دو مؤلفه سرعت نسبی‌اند و مطابق روابط زیر نرمال‌سازی می‌شوند:

\begin{equation}
	\tilde{x}_i=\mathrm{clip}\left(\frac{x_i}{15},-1,1\right)\quad,
	\tilde{y}_i=\mathrm{clip}\left(\frac{y_i}{15},-1,1\right)\quad,
	\tilde{d}_i=\mathrm{clip}\left(\frac{d_i}{10},0,1\right)
	\label{eq:obstacle_position_normalization_final}
\end{equation}

\begin{equation}
	\tilde{v}^{rel}_{x,i}=\mathrm{clip}\left(\frac{v^{rel}_{x,i}}{2},-1,1\right)\quad,
	\tilde{v}^{rel}_{y,i}=\mathrm{clip}\left(\frac{v^{rel}_{y,i}}{2},-1,1\right)
	\label{eq:obstacle_velocity_normalization_final}
\end{equation}

موانع از نزدیک به دور مرتب می‌شوند. در مسیر دقیق، موقعیت از مدل‌های شبیه‌سازی و در مسیر ادراکی از گره تشخیص و عمق دریافت می‌شود. برای کاهش تأخیر، موقعیت از \lr{raw\_x\_m} و \lr{raw\_y\_m} و سرعت از خروجی هموار کالمن گرفته می‌شود. اسلات‌های خالی در فضای فیزیکی با \((0,0,d_{max},0,0)\) و پس از نرمال‌سازی با \([0,0,1,0,0]\) پر می‌شوند تا با مانع بسیار نزدیک اشتباه نشوند. شکل \ref{fig:observation_space_infographic} ساختار مشاهده را نشان می‌دهد.
\begin{figure}[ht]
	\centering
	\includegraphics[width=0.95\textwidth]{observation_space_infographic}
	\caption{ساختار فضای مشاهده عامل یادگیری تقویتی شامل بردار ویژگی‌های ثابت و ماتریس ویژگی‌های موانع}
	\label{fig:observation_space_infographic}
	
	\vspace{0.4em}
	
	{\small
		\href{https://mohammadabedi1179.github.io/thesis-archive/chapters/ch4/observation_space_infographic.html}
		{\textbf{مشاهده نسخه تعاملی شکل با کیفیت بالاتر}}
		\qquad
		\qrcode[height=1.2cm]{https://mohammadabedi1179.github.io/thesis-archive/chapters/ch4/observation_space_infographic.html}
	}
\end{figure}


\subsection{فضای عمل و اعمال فرمان}
\label{subsec:action_space_final}

فضای عمل پیوسته و دوبعدی و خروجی سیاست در بازه \([-1,1]^2\) است:

\begin{equation}
	\mathbf{a}_t =
	\begin{bmatrix}
		a_{v,t} & a_{\omega,t}
	\end{bmatrix}^{T}
	\qquad
	\mathbf{a}_t \in [-1,1]^2
	\label{eq:normalized_action_final}
\end{equation}

عمل نرمال‌شده به سرعت خطی و زاویه‌ای مطلوب تبدیل می‌شود:

\begin{equation}
	v^{des}_t = a_{v,t} v_{max}
	\qquad
	\omega^{des}_t = a_{\omega,t} \omega_{max}
	\label{eq:action_scaling_final}
\end{equation}

با \(v_{max}=1~m/s\) و \(\omega_{max}=1~rad/s\)، حرکت جلو و عقب ممکن است. پیش از اجرا، فرمان مطلوب با قصد قبلی سیاست هموار می‌شود:

\begin{equation}
	v^{policy}_t = (1-\alpha_s)v^{policy}_{t-1} + \alpha_s v^{des}_t
	\qquad
	\omega^{policy}_t = (1-\alpha_s)\omega^{policy}_{t-1} + \alpha_s \omega^{des}_t
	\label{eq:policy_command_smoothing_final}
\end{equation}

هموارسازی بر فرمان قبلی سیاست، نه فرمان احتمالی اصلاح‌شده سپر، انجام می‌شود تا اثر مداخله پس از رفع خطر باقی نماند. فرمان حاصل برای ارزیابی خطر به سپر داده می‌شود؛ در آموزش، سپر نظارتی است و فرمان را تغییر نمی‌دهد، اما سطح خطر در مشاهده و پاداش وارد می‌شود. اگر فرمان اجرایی چند گام نزدیک صفر بماند، منطق مستقل فرار از بن‌بست چرخشی اجباری به سمت فضای بازتر اعمال می‌کند.

%----------------------------------------------------------
%----------------------------------------------------------
\subsection{سپر ایمنی و تعیین سطح خطر}
\label{subsec:safety_shield_levels_final}

در این پژوهش، سپر ایمنی به‌عنوان یک لایه تشخیص سطح خطر در محیط تعریف شده است. وظیفه این لایه، ارزیابی وضعیت لحظه‌ای ربات نسبت به موانع نزدیک و تعیین شدت خطر بر اساس داده‌های حسگرهای فراصوت و موانع مشاهده‌شده است. خروجی سپر ایمنی یکی از چهار سطح زیر است:

\begin{equation}
	\ell_t \in 
	\{
	\text{ایمن},
	\text{احتیاط},
	\text{خطرناک},
	\text{بحرانی}
	\}
	\label{eq:shield_levels_final}
\end{equation}

منطق تعیین سطح سپر به‌صورت سلسله‌مراتبی تعریف شده است؛ به این معنا که شرایط سطح شدیدتر اولویت بالاتری دارند. بنابراین، اگر وضعیت بحرانی تشخیص داده شود، همان سطح به‌عنوان خروجی سپر انتخاب می‌شود و سطوح ضعیف‌تر مانند خطرناک یا احتیاط جایگزین آن نخواهند شد. این طراحی باعث می‌شود سطح خطر گزارش‌شده همواره بیانگر نزدیک‌ترین و جدی‌ترین وضعیت ایمنی در گام زمانی جاری باشد.

در پیاده‌سازی انجام‌شده، سپر ایمنی دو منبع اطلاعاتی اصلی را بررسی می‌کند. منبع نخست، فاصله‌های نزدیک‌برد حاصل از حسگرهای فراصوت است که برای تشخیص خطرهای بسیار نزدیک و نواحی کور حسگر دیداری استفاده می‌شود. منبع دوم، موقعیت موانع ادراکی یا هندسی در ناحیه جلویی ربات است. برای این منظور، تنها موانعی در منطق سپر جلو در نظر گرفته می‌شوند که در مقابل ربات قرار داشته و از نظر جانبی در محدوده کریدور حرکتی ربات باشند.برای روشن‌تر شدن هندسه نواحی خطر، شکل \ref{fig:safety_shield_zones} نمای شماتیک دو منبع اصلی مورد استفاده در سپر ایمنی را نشان می‌دهد. شکل الف نواحی نزدیک‌برد حاصل از حسگرهای فراصوت را نمایش می‌دهد که برای تشخیص خطرهای بسیار نزدیک در اطراف بدنه استفاده می‌شوند. در مقابل، شکل ب ناحیه جلویی ربات را بر اساس فاصله مانع مشاهده‌شده نشان می‌دهد؛ این ناحیه برای ارزیابی موانعی به کار می‌رود که در مسیر حرکت ربات و در محدوده کریدور حرکتی آن قرار دارند.

\begin{figure}[htbp]
	\centering
	\setlength{\tabcolsep}{6pt}
	
	\begin{tabular}{cc}
		\includegraphics[width=0.47\textwidth]{ultrasonic_safety_shield.png}
		&
		\includegraphics[width=0.47\textwidth]{front_obstacle_shield.png}
		\\[4pt]
		\small الف) نواحی خطر مبتنی بر حسگرهای فراصوت
		&
		\small ب) نواحی خطر مربوط به مانع مشاهده‌شده در جلوی ربات
	\end{tabular}
	
	\caption{نمای شماتیک نواحی سپر ایمنی بر اساس دو منبع اطلاعاتی مورد استفاده در این پژوهش}
	\label{fig:safety_shield_zones}
	\vspace{0.4em}
	{\small
		\href{https://mohammadabedi1179.github.io/thesis-archive/chapters/ch4/ultrasonic-safety-shield.html}
		{\textbf{مشاهده نسخه تعاملی شکل با کیفیت بالاتر}}
		\qquad
		\qrcode[height=1.2cm]{https://mohammadabedi1179.github.io/thesis-archive/chapters/ch4/ultrasonic-safety-shield.html}
	}
\end{figure}

در جدول \ref{tab:shield_thresholds_final} آستانه‌های مورد استفاده برای تعیین سطح سپر نشان داده شده است. در این جدول، \(d_{front}\) کمترین فاصله دو حسگر فراصوت جلویی، \(d_{side}\) کمترین فاصله حسگرهای فراصوت کناری، و \(d^{obs}_{front}\) فاصله نزدیک‌ترین مانع مشاهده‌شده در ناحیه جلویی ربات است.

\begin{table}[ht]
	\centering
	\caption{آستانه‌های تعیین سطح سپر ایمنی}
	\label{tab:shield_thresholds_final}
	\resizebox{\textwidth}{!}{%
		\begin{tabular}{|c|c|c|c|}
			\hline
			\textbf{سطح سپر} & \textbf{شرط مبتنی بر فراصوت} & \textbf{شرط مبتنی بر مانع جلویی} & \textbf{شدت در مشاهده} \\
			\hline
			ایمن 
			& خارج از نواحی خطر 
			& خارج از نواحی خطر 
			& \(0.0\) \\
			\hline
			احتیاط 
			& \(d_{front}<0.80~m\) یا \(d_{side}<0.50~m\) 
			& \(d^{obs}_{front}<2.5~m\) 
			& \(0.33\) \\
			\hline
			خطرناک 
			& \(d_{front}<0.50~m\) 
			& \(d^{obs}_{front}<1.5~m\) 
			& \(0.67\) \\
			\hline
			بحرانی 
			& \(\min(d_{front},d_{side}, d_{rear})<0.25~m\) 
			& \(d^{obs}_{front}<0.9~m\) 
			& \(1.0\) \\
			\hline
		\end{tabular}%
	}
\end{table}

در حالت کلی، منطق سپر ایمنی قابلیت مداخله مستقیم در فرمان حرکتی را نیز دارد؛ برای مثال، در سطح بحرانی می‌تواند سرعت خطی را کاهش دهد، فرمان چرخش ایمن ایجاد کند یا در صورت باز بودن فضای عقب، فرمان عقب‌روی محدود اعمال نماید. با این حال، در پیکربندی آموزشی نهایی این پژوهش، سپر در حالت غیرفعال از نظر مداخله مستقیم، استفاده شده است. بنابراین، فرمان تولیدشده توسط سیاست یادگیری تقویتی در سطوح مختلف سپر بازنویسی نمی‌شود و مستقیماً به لایه کنترل حرکتی ارسال می‌گردد.

در این حالت، نقش سپر ایمنی به سه بخش محدود می‌شود. نخست، سطح خطر محیط را در هر گام زمانی تشخیص می‌دهد. دوم، سطح خطر و مدت حضور پیوسته در وضعیت سپر را به‌صورت دو ویژگی \(\ell_{shield}\) و \(\tau_{shield}\) وارد بردار مشاهده عامل می‌کند. سوم، متناسب با شدت سطح خطر، یک هزینه آموزشی مرحله‌ای تولید می‌کند که در تابع پاداش استفاده می‌شود. به این ترتیب، عامل به‌جای وابستگی مستقیم به یک کنترل‌کننده مداخله‌گر، در طول آموزش یاد می‌گیرد از ورود به نواحی خطر اجتناب کند.

نگاشت عددی سطح سپر در بردار مشاهده به‌صورت زیر تعریف شده است:

\begin{equation}
	\ell_{shield} =
	\begin{cases}
		0.0 & \ell_t=\text{ایمن},\\
		0.33 & \ell_t=\text{احتیاط},\\
		0.67 & \ell_t=\text{خطرناک},\\
		1.0 & \ell_t=\text{بحرانی}.
	\end{cases}
	\label{eq:shield_level_encoding_final}
\end{equation}

همچنین مدت حضور پیوسته در وضعیت سپر به‌صورت نرمال‌شده وارد مشاهده می‌شود:

\begin{equation}
	\tau_{shield} =
	\min
	\left(
	\frac{N_{shield}}{50},
	1
	\right)
	\label{eq:shield_duration_encoding_final}
\end{equation}

که در آن \(N_{shield}\) تعداد گام‌های متوالی حضور عامل در یکی از سطوح غیرایمن سپر است. این ویژگی به عامل کمک می‌کند میان ورود کوتاه‌مدت و گذرای ربات به ناحیه خطر و باقی‌ماندن طولانی‌تر در وضعیت نامطلوب تمایز قائل شود.

در ادامه، سطح سپر به یک هزینه مرحله‌ای تبدیل شده و در تابع پاداش مورد استفاده قرار می‌گیرد. جزئیات مربوط به هزینه سپر، ضریب لاگرانژ و تنظیم خودکار شدت جریمه ایمنی در بخش \ref{subsubsec:shield_lagrangian_reward_final} ارائه شده است.

\subsection{طراحی تابع پاداش و شرایط پایان اپیزود}
\label{subsec:reward_design_final}

تابع پاداش در این پژوهش به‌صورت ترکیبی طراحی شده است تا هم‌زمان چند هدف اصلی را پوشش دهد: حرکت به سمت هدف، کاهش زمان رسیدن، اجتناب از موانع، کاهش فرمان‌های نوسانی، حفظ فاصله ایمن و دوری از نواحی خطر. در هر گام زمانی، پاداش کلی به‌صورت زیر محاسبه می‌شود:

\begin{equation}
	r_t =
	-r_{step}
	+r_{TTC}
	+r_{progress}
	+r_{heading}
	+r_{action}
	+r_{ultra}
	+r_{shield}
	+r_{terminal}
	\label{eq:total_reward_final}
\end{equation}

در این رابطه، \(r_{step}=0.2\) جریمه ثابت گام زمانی است که در هر گام از پاداش کسر می‌شود. این جریمه عامل را به اتمام سریع‌تر مأموریت تشویق می‌کند.

%-----------------------------
\subsubsection{جریمه مبتنی بر زمان تا برخورد}
\label{subsubsec:ttc_reward_final}

برای کاهش رفتارهای پرخطر در نزدیکی موانع جلویی، از معیار زمان تا برخورد«%
\LTRfootnote{Time to Collision-TTC}» %
 استفاده شده است. مقدار \lr{TTC} برای موانعی محاسبه می‌شود که در جلوی ربات قرار دارند و فاصله جانبی آن‌ها کمتر از \(0.8~m\) است. اگر \(x_i\) فاصله طولی مانع و \(v^{rel}_{x,i}\) سرعت نسبی آن در راستای طولی باشد، سرعت نزدیک‌شدن به‌صورت زیر محاسبه می‌شود:

\begin{equation}
	v^{app}_i = v^{cmd}_t - v^{rel}_{x,i}
	\label{eq:approach_speed_final}
\end{equation}

اگر \(v^{app}_i>0\)، زمان تا برخورد برای آن مانع برابر است با:

\begin{equation}
	TTC_i = \frac{x_i}{v^{app}_i}
	\label{eq:ttc_single_final}
\end{equation}

کمترین مقدار \(TTC_i\) در میان موانع معتبر به‌عنوان \(TTC\) گام جاری در نظر گرفته می‌شود. اگر سرعت اجرایی رو به جلو بزرگ‌تر از \(0.1~m/s\) باشد و \(TTC<2.0~s\)، جریمه زیر اعمال می‌شود:

\begin{equation}
	r_{TTC}=
	\begin{cases}
		-\exp(-TTC), & v^{cmd}_t>0.1 \;\; \text{و} \;\; TTC<2.0\\
		0, & \text{در غیر این صورت}
	\end{cases}
	\label{eq:ttc_reward_term_final}
\end{equation}

در اینجا \(v^{cmd}_t\) سرعت اجرایی پس از منطق فرمان‌دهی است، نه صرفاً خروجی خام شبکه سیاست. بنابراین عامل بر اساس فرمانی تنبیه می‌شود که واقعاً در محیط اعمال شده است.

%-----------------------------
\subsubsection{پاداش پیشرفت به سمت هدف}
\label{subsubsec:progress_reward_final}

پاداش پیشرفت بر اساس کاهش فاصله ربات تا هدف بین دو گام متوالی تعریف شده است:

\begin{equation}
	r_{progress} = c_{progress}(d_{t-1}-d_t)
	\qquad
	c_{progress}=2.0
	\label{eq:progress_reward_final}
\end{equation}

اگر ربات به هدف نزدیک‌تر شود، این جمله مثبت خواهد بود و اگر از هدف دور شود، مقدار منفی ایجاد می‌کند. این مؤلفه یکی از مهم‌ترین اجزای پاداش چگال است و باعث می‌شود عامل پیش از دریافت پاداش نهایی نیز بازخورد جهت‌دار داشته باشد.

%-----------------------------
\subsubsection{پاداش تطبیقی هم‌راستایی با هدف}
\label{subsubsec:adaptive_heading_reward_final}

در نسخه نهایی، پاداش هم‌راستایی صرفاً بر اساس \(\cos(\phi_g)\) محاسبه نمی‌شود. دلیل این تغییر آن است که پاداش کسینوسی می‌تواند حتی در شرایطی که مسیر مستقیم به سمت هدف توسط موانع مسدود شده است، عامل را به چرخیدن یا حرکت نامناسب تشویق کند. بنابراین، در این پژوهش یک پاداش تطبیقی استفاده شده است که ابتدا جهت‌های مسدودشده توسط موانع نزدیک را حذف کرده و سپس نزدیک‌ترین جهت آزاد به جهت هدف را پیدا می‌کند.

برای هر مانع نزدیک‌تر از \(5.0~m\) که در نیم‌صفحه جلویی ربات قرار داشته باشد، یک مخروط زاویه‌ای مسدود در اطراف زاویه مانع تعریف می‌شود. نیم‌عرض پایه این مخروط \(0.25~rad\) است و برای موانع نزدیک‌تر بزرگ‌تر در نظر گرفته می‌شود:

\begin{equation}
	\Delta\theta^{obs}_i
	=0.25\left(1+\max\left(0,1-\frac{d_i}{5}\right)\right)
	\label{eq:obstacle_blocking_cone_final}
\end{equation}

زاویه هدف ابتدا به محدوده میدان دید جلویی \([-\pi/2,\pi/2]\) محدود می‌شود. سپس با گام زاویه‌ای \(0.1~rad\)، جهت‌های اطراف زاویه هدف جست‌وجو می‌شوند تا اولین جهت آزاد \(\theta^*\) پیدا شود. اگر جهت آزادی پیدا نشود، اگر ربات در حال عقب‌روی باشد، یا اگر سطح سپر بحرانی باشد، پاداش هم‌راستایی صفر می‌شود. در غیر این صورت، پاداش به‌صورت زیر محاسبه می‌گردد:

\begin{equation}
	r_{heading}=
	\begin{cases}
		c_h\left(\theta_{max}-|\theta^*-\phi_g|\right)
		\mathrm{clip}\left(\frac{v^{cmd}_t}{v_{max}},0,1\right),
		& |\theta^*-\phi_g|\le\theta_{max}\\
		0, & \text{در غیر این صورت}
	\end{cases}
	\label{eq:adaptive_heading_reward_final}
\end{equation}

که در آن \(c_h=0.6\) و \(\theta_{max}=\pi/6\) است. این فرمول عامل را تشویق می‌کند در شرایط ایمن به سمت هدف حرکت کند، اما در حضور مانع، جهت آزاد نزدیک به هدف را ترجیح دهد.

%-----------------------------
\subsubsection{جریمه تغییرات شدید فرمان}
\label{subsubsec:action_smooth_reward_final}

برای کاهش رفتارهای نوسانی و تولید فرمان‌های نرم‌تر، تغییرات فرمان سیاست بین دو گام متوالی جریمه می‌شود:

\begin{equation}
	r_{action} = -c_{action}\left(|\Delta v|+|\Delta\omega|\right)
	\qquad
	c_{action}=0.2
	\label{eq:action_change_penalty_final}
\end{equation}

در این رابطه، \(\Delta v\) و \(\Delta\omega\) بر اساس فرمان هموارشده سیاست محاسبه می‌شوند. این جریمه کمک می‌کند عامل از تغییرات شدید سرعت خطی و زاویه‌ای پرهیز کند.

%-----------------------------
\subsubsection{جریمه نزدیکی بر اساس حسگرهای فراصوت}
\label{subsubsec:ultrasonic_reward_final}

برای افزایش ایمنی نزدیک‌برد، اگر کمترین فاصله حسگرهای فراصوت جلو و کناری کمتر از \(0.5~m\) شود، جریمه زیر اعمال می‌گردد:

\begin{equation}
	r_{ultra}=
	\begin{cases}
		-2(0.5-d_{min}^{ultra}), & d_{min}^{ultra}<0.5\\
		0, & \text{در غیر این صورت}
	\end{cases}
	\label{eq:ultra_reward_final}
\end{equation}

در اینجا \(d_{min}^{ultra}\) کمترین فاصله میان چهار حسگر جلو-چپ، جلو-راست، سمت چپ، سمت راست و عقب است.
%----------------------------------------------------------
\subsubsection{هزینه سپر ایمنی و تنظیم لاگرانژی}
\label{subsubsec:shield_lagrangian_reward_final}

برای تبدیل سطح سپر ایمنی به بازخورد آموزشی، یک هزینه مرحله‌ای شدت‌دار تعریف شده است. این هزینه با افزایش شدت خطر بیشتر می‌شود و به عامل یادگیری تقویتی نشان می‌دهد که ورود به نواحی نزدیک‌تر به برخورد نامطلوب‌تر است. هزینه مرحله‌ای سپر به‌صورت زیر تعریف می‌شود:

\begin{equation}
	c_t =
	\begin{cases}
		0.0 & \ell_t=\text{ایمن}\\
		0.4 & \ell_t=\text{احتیاط}\\
		1.5 & \ell_t=\text{خطرناک}\\
		3.0 & \ell_t=\text{بحرانی}
	\end{cases}
	\label{eq:shield_step_cost_final}
\end{equation}

بنابراین، برخلاف جریمه برخورد که تنها در زمان وقوع برخورد اعمال می‌شود، هزینه سپر یک بازخورد مرحله‌ای و پیوسته‌تر فراهم می‌کند. این طراحی باعث می‌شود عامل پیش از رسیدن به وضعیت برخورد، نسبت به نزدیک‌شدن به مانع بازخورد منفی دریافت کند و به‌تدریج سیاست ایمن‌تری بیاموزد.

هزینه تجمعی سپر در هر اپیزود از مجموع هزینه‌های مرحله‌ای به‌دست می‌آید:

\begin{equation}
	J_c =
	\sum_{t=0}^{T} c_t
	\label{eq:shield_episode_cost_final}
\end{equation}

برای جلوگیری از آن‌که مقدار تجمعی هزینه سپر از نظر عددی بیش از حد بزرگ شود و بر سایر اجزای پاداش، به‌ویژه پاداش رسیدن به هدف، غلبه کند، این هزینه در محاسبات آموزشی با سقف مشخصی محدود می‌شود:

\begin{equation}
	J_c \leftarrow \min(J_c,25)
	\label{eq:shield_cost_cap_final}
\end{equation}

جمله مربوط به سپر در تابع پاداش به‌صورت زیر وارد می‌شود:

\begin{equation}
	r_{shield,t}
	=
	-\lambda_t c_t
	\label{eq:shield_reward_term_final}
\end{equation}

در این رابطه، \(\lambda_t\) ضریب لاگرانژ است که شدت اثر هزینه سپر در پاداش را تعیین می‌کند. اگر \(\lambda_t\) بسیار کوچک باشد، عامل ممکن است برای رسیدن سریع‌تر به هدف، ورود به نواحی خطر را نادیده بگیرد. در مقابل، اگر \(\lambda_t\) بیش از حد بزرگ انتخاب شود، عامل ممکن است بیش از اندازه محافظه‌کار شود و یادگیری رسیدن به هدف مختل گردد. بنابراین، انتخاب یک مقدار ثابت و دستی برای این ضریب می‌تواند باعث رفتار نامتعادل در فرایند آموزش شود.

برای رفع این مشکل، در این پژوهش مقدار \(\lambda_t\) به‌صورت تطبیقی و بر اساس میزان نقض قید ایمنی در اپیزودهای آموزشی تنظیم می‌شود. ابتدا هزینه اپیزودی سپر با استفاده از یک میانگین نمایی هموار می‌شود:

\begin{equation}
	\hat{J}^{(e)}_c =
	(1-\beta)\hat{J}^{(e-1)}_c
	+
	\beta J^{(e)}_c
	\qquad
	\beta=0.1
	\label{eq:shield_cost_ema_final}
\end{equation}

در این رابطه، \(e\) شماره اپیزود آموزشی و \(\hat{J}^{(e)}_c\) هزینه هموارشده سپر تا اپیزود \(e\) است. سپس خطای قید ایمنی به‌صورت اختلاف میان هزینه هموارشده و بودجه مجاز تعریف می‌شود:

\begin{equation}
	e_c^{(e)} =
	\hat{J}^{(e)}_c
	-
	J_{budget}
	\qquad
	J_{budget}=6.0
	\label{eq:shield_constraint_error_final}
\end{equation}

اگر \(e_c^{(e)}>0\) باشد، هزینه سپر از بودجه تعیین‌شده بیشتر شده و لازم است شدت جریمه ایمنی افزایش یابد. در مقابل، اگر هزینه سپر کمتر از بودجه باشد، نیازی به افزایش ضریب لاگرانژ وجود ندارد. برای تنظیم پایدار این ضریب، از یک ساختار \lr{PID} استفاده شده است:

\begin{equation}
	\lambda^{*} =
	\mathrm{clip}
	\left(
	K_P e_c^{(e)}
	+
	K_I I^{(e)}
	+
	K_D D^{(e)},
	0,
	2
	\right).
	\label{eq:pid_lambda_raw_final}
\end{equation}

در این رابطه، ضرایب کنترل‌کننده برابر با

\begin{equation}
	K_P=0.05
	\qquad
	K_I=0.01
	\qquad
	K_D=0.02
	\label{eq:pid_gains_final}
\end{equation}

در نظر گرفته شده‌اند. جمله انتگرالی و مشتقی نیز به‌صورت زیر تعریف می‌شوند:

\begin{equation}
	I^{(e)}
	=
	\max
	\left(
	0,
	I^{(e-1)}+e_c^{(e)}
	\right)
	\label{eq:pid_integral_final}
\end{equation}

\begin{equation}
	D^{(e)}
	=
	e_c^{(e)}
	-
	e_c^{(e-1)}
	\label{eq:pid_derivative_final}
\end{equation}

استفاده از ساختار \lr{PID} در این بخش به‌منظور کنترل مستقیم حرکت ربات نیست، بلکه برای تنظیم تطبیقی ضریب لاگرانژ در سطح آموزش استفاده شده است. جزء تناسبی به خطای فعلی هزینه سپر نسبت به بودجه واکنش نشان می‌دهد، جزء انتگرالی اثر نقض‌های تکرارشونده قید ایمنی را در چندین اپیزود متوالی لحاظ می‌کند و جزء مشتقی نسبت به روند تغییرات خطا حساس است. به این ترتیب، شدت جریمه ایمنی متناسب با رفتار عامل در طول آموزش تنظیم می‌شود و نیازی به انتخاب دستی یک وزن ثابت برای تمام مراحل آموزش وجود ندارد.

برای جلوگیری از تداخل زودهنگام جریمه ایمنی با یادگیری اولیه عامل، ضریب لاگرانژ در ابتدای آموزش به‌صورت تدریجی فعال می‌شود. در \(150\) اپیزود نخست، مقدار \(\lambda\) برابر صفر در نظر گرفته می‌شود تا عامل ابتدا رفتار پایه رسیدن به هدف را یاد بگیرد. سپس طی \(150\) اپیزود بعدی، اثر ضریب لاگرانژ به‌صورت تدریجی افزایش می‌یابد:

\begin{equation}
	\lambda^{(e)}
	=
	\alpha^{(e)} \lambda^{*}
	\label{eq:lambda_ramp_final}
\end{equation}

که در آن ضریب فعال‌سازی تدریجی به‌صورت زیر تعریف می‌شود:

\begin{equation}
	\alpha^{(e)}
	=
	\begin{cases}
		0,
		& e < 150\\
		\dfrac{e-150}{150},
		& 150 \le e < 300\\
		1,
		& e \ge 300
	\end{cases}
	\label{eq:lambda_warmup_ramp_final}
\end{equation}

در نتیجه، جمله سپر ایمنی در تابع پاداش نهایی به‌صورت یک جریمه تطبیقی عمل می‌کند؛ به‌گونه‌ای که اگر عامل به‌طور مکرر وارد نواحی خطر شود، ضریب لاگرانژ افزایش یافته و هزینه سپر اثر بیشتری بر یادگیری خواهد داشت. در مقابل، اگر عامل رفتار ایمن‌تری نشان دهد، شدت جریمه در محدوده کنترل‌شده باقی می‌ماند و از محافظه‌کاری بیش از حد جلوگیری می‌شود.

%-----------------------------
\subsubsection{پاداش‌ها و جریمه‌های نهایی}
\label{subsubsec:terminal_reward_final}

سه حالت پایانی اصلی در محیط وجود دارد: برخورد، رسیدن به هدف و پایان زمان. اگر کمترین فاصله فراصوتی جلو/کناری کمتر از \(0.2~m\) شود یا فاصله مانع ادراکی/هندسی کمتر از \(0.7~m\) گردد، اپیزود با برخورد خاتمه می‌یابد و جریمه \(-100\) اعمال می‌شود:

\begin{equation}
	r_{collision}=-100
	\label{eq:collision_penalty_final}
\end{equation}


در این پژوهش، تحقق هدف صرفاً به کاهش فاصله ربات تا موقعیت هدف محدود نمی‌شود، بلکه یک شرط چندعاملی و هم‌زمان تعریف شده است که برای جلوگیری از تشخیص‌های نادرست و توقف‌های زودهنگام در محیط‌های دارای مانع ضروری است.

به‌طور مشخص، وضعیت «رسیدن به هدف» تنها در صورتی معتبر است که هر سه شرط زیر به‌صورت هم‌زمان برقرار باشند:

\begin{enumerate}
	\item فاصله خطی ربات تا هدف کمتر از شعاع موفقیت تعریف‌شده باشد؛
	\item قدرمطلق زاویه انحراف هدف نسبت به جهت حرکت ربات کمتر از \(\pi/6\) باشد؛
	\item هیچ مانعی خط دید مستقیم میان ربات و هدف را مسدود نکند.
\end{enumerate}

برای شرط سوم، فاصله هر مانع تا پاره‌خط میان ربات و هدف محاسبه می‌شود. اگر مانعی با شعاع مؤثر \(0.3~m\) بین ربات و هدف قرار گیرد، هدف هنوز محقق‌شده در نظر گرفته نمی‌شود. در صورت تحقق موفق هدف، پاداش نهایی زیر اعمال می‌شود:

\begin{equation}
	r_{goal}=+100
	\label{eq:goal_reward_final}
\end{equation}

این تعریف در شکل \ref{fig:goal_conditions} به‌صورت تصویری نمایش داده شده است. همان‌طور که مشاهده می‌شود، تنها در حالتی که هر سه شرط به‌صورت هم‌زمان برقرار باشند، سیستم وضعیت رسیدن به هدف را معتبر تلقی می‌کند و در غیر این صورت، حتی در صورت برقراری دو شرط از سه شرط، اپیزود موفق محسوب نخواهد شد.

%----------------------------------------------------------
\begin{figure}[ht]

\end{figure}

\begin{figure}[ht]
	\centering
\includegraphics[width=0.95\textwidth]{goal_conditions.png}
\caption{شرایط تحقق هدف در سامانه ناوبری پیشنهادی}
\label{fig:goal_conditions}
	
	\vspace{0.4em}
	
	{\small
		\href{https://mohammadabedi1179.github.io/thesis-archive/chapters/ch4/goal-conditions.html}
		{\textbf{مشاهده نسخه تعاملی شکل با کیفیت بالاتر}}
		\qquad
		\qrcode[height=1.2cm]{https://mohammadabedi1179.github.io/thesis-archive/chapters/ch4/goal-conditions.html}
	}
\end{figure}
%----------------------------------------------------------

اگر اپیزود بدون برخورد و بدون رسیدن به هدف از حد زمانی مرحله فراتر رود، اپیزود خاتمه یافته و جریمه \(-10\) اعمال می‌شود:

\begin{equation}
	r_{timeout}=-10
	\label{eq:timeout_penalty_final}
\end{equation}

%----------------------------------------------------------
\subsection{معماری سیاست \lr{LSTM-SAC}}
\label{subsec:lstm_sac_architecture_final}

عامل از \lr{SAC} با استخراج‌گر سفارشی \lr{LSTM} استفاده می‌کند. \lr{LSTM} در اینجا حافظه پنهان میان گام‌های اپیزود را نگه نمی‌دارد، بلکه دنباله موانع مرتب‌شده همان مشاهده را فشرده می‌کند؛ اطلاعات زمانی از سرعت‌های نسبی، حافظه محیط و کالمن وارد شده‌اند. بردار ثابت \(D_f=15\) مستقیم و ماتریس \(10\times5\) به \lr{LSTM} داده می‌شود. برای دسته \(B\):

\begin{equation}
	\mathbf{X}^{obs}\in\mathbb{R}^{B\times 10\times 5}
	\label{eq:lstm_input_shape_final}
\end{equation}

شبکه یک لایه با حالت پنهان \(H=64\) دارد:

\begin{equation}
	\mathrm{LSTM}(input\_size=5,\;hidden\_size=64,\;num\_layers=1)
	\label{eq:lstm_config_final}
\end{equation}

حالت نهایی، نمایش فشرده موانع است:

\begin{equation}
	\mathbf{h}^{obs}=\mathbf{H}_{n}[-1]\in\mathbb{R}^{B\times64}
	\label{eq:lstm_last_hidden_final}
\end{equation}

این نمایش با بردار ثابت الحاق می‌شود:

\begin{equation}
	\mathbf{z}_t=\left[\mathbf{s}^{fixed}_t\;\|\;\mathbf{h}^{obs}_t\right]
	\label{eq:feature_concat_final}
\end{equation}

بعد بردار ترکیبی در پیکربندی حاضر برابر است با:

\begin{equation}
	\dim(\mathbf{z}_t)=15+64=79
	\label{eq:combined_feature_dim_final}
\end{equation}

بعد ورودی از فضای مشاهده به‌صورت پویا محاسبه می‌شود تا تغییر ابعاد ثابت یا حالت پنهان خودکار اعمال گردد. شکل \ref{fig:lstm_sac_policy_architecture} دو مسیر موازی بردار ثابت و ماتریس موانع و الحاق آن‌ها در \(\mathbf{z}_t\) را نشان می‌دهد.

\begin{figure}[htbp]
\end{figure}

\begin{figure}[htbp]
	\centering
\includegraphics[width=0.75\textwidth]{LSTM_SAC_architecture}
\caption{معماری سیاست \lr{LSTM-SAC} در عامل پیشنهادی.}
\label{fig:lstm_sac_policy_architecture}
	\vspace{0.4em}
	{\small
		\href{https://mohammadabedi1179.github.io/thesis-archive/chapters/ch4/LSTM_SAC_architecture.html}
		{\textbf{مشاهده نسخه تعاملی شکل با کیفیت بالاتر}}
		\qquad
		\qrcode[height=1.2cm]{https://mohammadabedi1179.github.io/thesis-archive/chapters/ch4/LSTM_SAC_architecture.html}
	}
\end{figure}

مطابق شکل \ref{fig:lstm_sac_policy_architecture}، کنشگر از \(\mathbf{z}_t\) پارامترهای توزیع و با \lr{Reparameterization} عمل \(\mathbf{a}_t=(v,\omega)\) را تولید می‌کند. دو منتقد مستقل زوج \((\mathbf{z}_t,\mathbf{a}_t)\) را برای کاهش بیش‌برآورد \(Q\) ارزیابی می‌کنند. شبکه‌های کنشگر و منتقد دارای لایه‌های \([256,256]\) و فعال‌سازی \lr{ReLU} هستند و عمل به‌صورت تصادفی یا قطعی استخراج می‌شود.
%----------------------------------------------------------
\subsection{آموزش معلم، دانشجو و لنگر رفتاری}
\label{subsec:teacher_student_distillation_final}

آموزش در دو حالت انجام می‌شود. معلم با \lr{SAC} و هندسه دقیق، سیاستی پایدار در مشاهده تمیز می‌آموزد. دانشجو با \lr{DistilledSAC} روی داده ادراکی آموزش می‌بیند و زیان تقلیدی را به هدف \lr{SAC} می‌افزاید. وزن‌های دانشجو در آغاز از معلم مقداردهی می‌شوند و در هر گام، مشاهده دقیق برای «معلم منجمد‌شده%
\LTRfootnote{Frozen Teacher}» %
ساخته می‌شود تا عمل قطعی زیر تولید گردد:
\begin{equation}
	a^{teacher}_t = \pi^{det}_{teacher}\left(o^{GT}_t\right)
	\label{eq:teacher_action_final}
\end{equation}

مشاهده دقیق فقط هدف رفتاری معلم را می‌سازد و گذار دانشجو با مشاهده ادراکی در حافظه بازپخش ذخیره می‌شود. حافظه سفارشی علاوه بر گذار، \(a^{teacher}\) و ماسک اعتبار \(m\) را نگه می‌دارد. زیان کنشگر عبارت است از:

\begin{equation}
	\mathcal{L}_{actor}=
	\mathcal{L}_{SAC}+\lambda_{BC}\mathcal{L}_{BC}
	\label{eq:distilled_actor_loss_final}
\end{equation}

زیان تقلیدی از فاصله عمل قطعی دانشجو و معلم به‌دست می‌آید:

\begin{equation}
	\mathcal{L}_{BC}=
	\frac{1}{N_{valid}D_a}
	\sum_{i=1}^{N}
	m_i
	\left\|
	\pi^{det}_{student}(o_i)-a^{teacher}_i
	\right\|_2^2
	\label{eq:bc_loss_final}
\end{equation}

\(m_i\) ماسک اعتبار، \(N_{valid}\) تعداد نمونه معتبر و \(D_a=2\) است. زیان تقلیدی تنها با نسبت اعتبار بیشتر از \(0.05\) فعال می‌شود. وزن \(\lambda_{BC}\) بر اساس مرحله تنظیم می‌شود؛ در انتقال ادراکی بزرگ‌تر و در مراحل استقلال دانشجو کوچک‌تر است (جدول \ref{tab:bc_weight_by_stage_final}).

\begin{table}[htbp]
	\centering
	\caption{وزن زیان تقلیدی در مراحل مختلف آموزش دانشجو}
	\label{tab:bc_weight_by_stage_final}
	\begin{tabular}{|c|c|}
		\hline
		\textbf{مرحله} & \(\lambda_{BC}\) \\
		\hline
		$0$ & $0.5$ \\
		\hline
		$1$ & $1.0$ \\
		\hline
		$2$ & $1.5$ \\
		\hline
		$3$ & $1.5$ \\
		\hline
		$4$ & $2.0$ \\
		\hline
		$5$ & $2.0$ \\
		\hline
		$6$ & $1.0$ \\
		\hline
		$7$ & $0.5$ \\
		\hline
		$8$ & $0.3$ \\
		\hline
		$9$ & $0.15$ \\
		\hline
	\end{tabular}
\end{table}

این زمان‌بندی، دانشجو را در انتقال دشوار نزدیک معلم نگه می‌دارد و در مراحل نهایی امکان اصلاح مستقل سیاست بر پایه ادراک را فراهم می‌کند.

%----------------------------------------------------------
\subsection{ابرپارامترهای آموزش}
\label{subsec:training_hyperparameters_final}

جدول \ref{tab:sac_hyperparams_final} ابرپارامترهای اصلی آموزش را نشان می‌دهد. این مقادیر مطابق اسکریپت آموزش نهایی تنظیم شده‌اند.

\begin{table}[htbp]
	\centering
	\caption{ابرپارامترهای اصلی آموزش \lr{SAC}/\lr{DistilledSAC}}
	\label{tab:sac_hyperparams_final}
	%\resizebox{\textwidth}{!}{%
		\begin{tabular}{|c|c|}
			\hline
			\textbf{پارامتر} & \textbf{مقدار} \\
			\hline
			الگوریتم معلم & \lr{SAC} \\
			\hline
			الگوریتم دانشجو & \lr{DistilledSAC} با حافظه بازپخش سفارشی \\
			\hline
			سیاست & \lr{LSTM-SAC} با \(H=64\) و یک لایه \lr{LSTM} \\
			\hline
			معماری \lr{MLP} کنشگر/منتقد & \([256,256]\) \\
			\hline
			اندازه حافظه بازپخش & \(10^6\) گذار \\
			\hline
			اندازه batch & \(256\) \\
			\hline
			ضریب تخفیف & \(\gamma=0.97\) \\
			\hline
			ضریب به‌روزرسانی نرم شبکه هدف & \(\tau=0.005\) \\
			\hline
			نرخ یادگیری اولیه & \(3\times 10^{-4}\) \\
			\hline
			زمان‌بندی نرخ یادگیری & \(lr=lr_0\max(0.5,progress)\) \\
			\hline
			شروع یادگیری & پس از \(5000\) گذار \\
			\hline
			نسبت به‌روزرسانی & یک گام گرادیانی به‌ازای هر گام محیط \\
			\hline
			تنظیم آنتروپی & خودکار\\
			\hline
			آنتروپی هدف & \(-1.0\) \\
			\hline
			تعداد کل گام‌های آموزش & \(3\times10^6\) \\
			\hline
			حالت سپر در آموزش & غیرفعال \\
			\hline
			بودجه هزینه سپر & \(J_{budget}=6.0\) \\
			\hline
			سقف ضریب لاگرانژ & \(\lambda_{max}=2.0\) \\
			\hline
		\end{tabular}%
		%}
\end{table}

%----------------------------------------------------------
\subsection{فرآیند آموزش عامل}
\label{subsec:training_process_final}

در هر گام، محیط دیکشنری مشاهده را می‌سازد، \lr{LSTM-SAC} عمل دوبعدی را تولید می‌کند، محیط آن را به فرمان هموار تبدیل کرده، سطح سپر را می‌سنجد و پس از اجرا، پاداش و مشاهده بعدی را بازمی‌گرداند. گذار معلم در حافظه عادی و گذار دانشجو همراه عمل معلم و ماسک اعتبار در حافظه سفارشی ذخیره می‌شود. در به‌روزرسانی، منتقدها، سپس کنشگر با زیان \lr{SAC} و در صورت اعتبار کافی زیان تقلیدی، و در پایان آنتروپی و شبکه‌های هدف اصلاح می‌شوند. الگوریتم \ref{alg:distilled_lstm_sac_training_final} این فرایند را خلاصه می‌کند.

\begin{algorithm}[htbp]
	\footnotesize
	\caption{فرآیند آموزش عامل \lr{LSTM-SAC} با سپر غیرفعال و لنگر رفتاری معلم}
	\label{alg:distilled_lstm_sac_training_final}
	\begin{algorithmic}[1]
		\REQUIRE محیط \(Env\)، سیاست \(\pi_\theta\)، منتقدها \(Q_{\phi_1},Q_{\phi_2}\)، حافظه \(\mathcal{D}\)، مدل معلم اختیاری \(\pi_T\)
		\ENSURE سیاست آموزش‌دیده \(\pi_\theta\)
		\STATE مقداردهی اولیه سیاست، منتقدها، شبکه‌های هدف و ضریب آنتروپی
		\FOR{هر گام تعامل با محیط}
		\STATE دریافت مشاهده \(o_t=\{\mathbf{s}^{fixed}_t,\mathbf{O}_t\}\)
		\STATE تولید عمل نرمال‌شده \(a_t\sim\pi_\theta(a|o_t)\)
		\IF{آموزش دانشجو و معلم موجود باشد}
		\STATE ساخت مشاهده موقت هندسه دقیق \(o^{GT}_t\) برای معلم
		\STATE محاسبه \(a^T_t=\pi_T^{det}(o^{GT}_t)\)
		\ENDIF
		\STATE تبدیل \(a_t\) به فرمان سرعت، اعمال هموارسازی، تعیین سطح سپر و اجرای فرمان
		\STATE دریافت \((o_{t+1},r_t,done,info_t)\)
		\STATE ذخیره گذار در \(\mathcal{D}\)؛ در حالت دانشجو، ذخیره \(a^T_t\) و ماسک اعتبار نیز انجام می‌شود
		\IF{تعداد گذارها بیشتر از \(5000\) باشد}
		\STATE نمونه‌گیری یک batch از \(\mathcal{D}\)
		\STATE محاسبه هدف منتقدها:
		\[
		y=r+\gamma(1-done)
		\left(
		\min_i Q_{\bar{\phi}_i}(o',a')-
		\alpha\log\pi_\theta(a'|o')
		\right)
		\]
		\STATE به‌روزرسانی \(Q_{\phi_1}\) و \(Q_{\phi_2}\) با کمینه‌سازی خطای مربعی
		\STATE محاسبه زیان کنشگر \(\mathcal{L}_{SAC}\)
		\IF{نمونه‌های معتبر معلم کافی باشد}
		\STATE محاسبه \(\mathcal{L}_{BC}\) و زیان کل \(\mathcal{L}_{actor}=\mathcal{L}_{SAC}+\lambda_{BC}\mathcal{L}_{BC}\)
		\ELSE
		\STATE استفاده از \(\mathcal{L}_{actor}=\mathcal{L}_{SAC}\)
		\ENDIF
		\STATE به‌روزرسانی کنشگر، ضریب آنتروپی و شبکه‌های هدف
		\ENDIF
		\IF{اپیزود تمام شده باشد}
		\STATE ثبت آمار اپیزود، به‌روزرسانی \(\lambda\) سپر و بررسی پیشرفت curriculum
		\ENDIF
		\ENDFOR
		\RETURN \(\pi_\theta\)
	\end{algorithmic}
\end{algorithm}

%----------------------------------------------------------
\subsection{سناریوهای آموزشی و یادگیری مرحله‌ای}
\label{subsec:curriculum_learning_final}

آموزش برای کاهش دشواری اولیه به‌صورت مرحله‌ای انجام می‌شود. مراحل صفر تا سه مهارت پایه را با مشاهده تمیز یا غالباً هندسی می‌آموزند؛ مراحل چهار تا شش با دشواری مأموریت تقریباً ثابت، انتقال از هندسه دقیق به ادراک را انجام می‌دهند؛ و از مرحله هفت، ادراک کامل با موانع و هدف متحرک به کار می‌رود. در حالت معلم، هندسه دقیق و آستانه عبور حداقل \(90\%\) حفظ می‌شود. در دانشجو، پارامتر ترکیب دستیار مشاهده، تعداد اپیزودهای کاهش تدریجی هندسه دقیق را تعیین می‌کند و مقدار صفر به معنای ادراک کامل است. مراحل در جدول \ref{tab:curriculum_stages_final} آمده‌اند.

\begin{table}[htbp]
	\centering
	\caption{مراحل یادگیری مرحله‌ای در آموزش عامل}
	\label{tab:curriculum_stages_final}
	\resizebox{\textwidth}{!}{%
		\begin{tabular}{|c|c|c|c|c|c|c|c|c|c|c|}
			\hline
			\textbf{مرحله} &
			\textbf{موانع ایستا} &
			\textbf{موانع متحرک} &
			\textbf{هدف متحرک} &
			\textbf{سرعت هدف} &
			\textbf{شعاع موفقیت} &
			\textbf{زمان} &
			\textbf{نیم‌زاویه قطاع} &
			\textbf{حداقل اپیزود} &
			\textbf{آستانه موفقیت} &
			\textbf{\lr{GT blend}} \\
			\hline
			0 & 0 & 0 & خیر & 0.0 & 1.5 & 80 & 15 & 100 & 90\% & 0 \\
			\hline
			1 & 2 & 0 & خیر & 0.0 & 1.5 & 80 & 20 & 120 & 88\% & 600 \\
			\hline
			2 & 1 & 1 & خیر & 0.0 & 1.5 & 80 & 20 & 150 & 85\% & 800 \\
			\hline
			3 & 2 & 2 & خیر & 0.0 & 1.5 & 85 & 30 & 150 & 80\% & 800 \\
			\hline
			4 & 2 & 0 & خیر & 0.0 & 1.5 & 85 & 30 & 180 & 82\% & 180 \\
			\hline
			5 & 2 & 2 & خیر & 0.0 & 1.5 & 85 & 30 & 250 & 78\% & 250 \\
			\hline
			6 & 2 & 2 & خیر & 0.0 & 1.5 & 85 & 30 & 200 & 80\% & 0 \\
			\hline
			7 & 2 & 3 & بله & 0.2 & 2.0 & 100 & 30 & 200 & 75\% & 0 \\
			\hline
			8 & 2 & 4 & بله & 0.5 & 2.0 & 100 & 45 & 250 & 72\% & 0 \\
			\hline
			9 & 3 & 4 & بله & 0.5 & 1.5 & 100 & 60 & 300 & 80\% & 0 \\
			\hline
		\end{tabular}%
	}
\end{table}

پس از حداقل اپیزود هر مرحله، نرخ موفقیت در پنجره لغزان ۵۰ تا ۲۰۰ اپیزودی بررسی می‌شود. با تحقق معیار، مدل ذخیره، نرخ یادگیری کاهش و تنظیمات مرحله بعد ارسال می‌گردد. ابتدای هر اپیزود، ربات بازنشانی و هدف روی دایره‌ای با شعاع تقریبی \(7~m\) و زاویه تصادفی قرار می‌گیرد؛ تعداد و قطاع موانع و در مراحل پایانی حرکت هدف نیز مطابق جدول تنظیم می‌شوند. عامل بدین‌ترتیب از رسیدن و اجتناب ساده به ادراک خطادار و محیط پویاتر منتقل می‌شود.

%----------------------------------------------------------
\subsection{جمع‌بندی طراحی عامل}
\label{subsec:agent_design_summary_final}

عامل از بردار ثابت ۱۵بعدی و ماتریس موانع \(10\times5\) استفاده می‌کند؛ \lr{LSTM} دنباله مرتب موانع همان گام را فشرده و اطلاعات زمانی از سرعت نسبی، حافظه و کالمن تأمین می‌شود. عمل دوبعدی به سرعت خطی و زاویه‌ای هموار تبدیل می‌گردد. پاداش، پیشرفت، \lr{TTC}، هم‌راستایی تطبیقی، نرمی، فراصوت، سپر و پایان اپیزود را ترکیب می‌کند. سپر غیرفعال با ویژگی مشاهده و هزینه لاگرانژی، و آموزش معلم--دانشجو با زیان مرحله‌محور، ایمنی و انتقال از هندسه دقیق به ادراک را تقویت می‌کنند.

%----------------------------------------------------------
\section{جمع‌بندی فصل چهارم}
\label{sec:chapter4_summary}

در این فصل، چارچوب پیشنهادی ناوبری خودران اسکوتر چهارچرخ در محیط‌های پویا از سطح مدل‌سازی تا آموزش عامل تشریح شد. معماری ماژولار سامانه، اجزای ادراک، تصمیم‌گیری، کنترل و ارزیابی را در یک حلقه ادراک--تصمیم--کنترل و در بستر \lr{ROS2}/\lr{Gazebo} یکپارچه کرد. مدل \lr{URDF/Xacro} نیز ساختار مکانیکی، جرم و اینرسی، چرخ‌ها، حسگرها و تماس با محیط را فراهم ساخت و لایه کنترل مبتنی بر \lr{skid-steering} فرمان سرعت خطی و زاویه‌ای سیاست را به سرعت چهار چرخ تبدیل نمود.

در لایه ادراک، تصاویر استریو برای تولید اختلاف مکانی و تخمین عمق به کار رفتند. آشکارسازی \lr{YOLO}، استخراج موقعیت دوبعدی موانع، حذف داده‌های پرت، پالایش کالمن و تخمین سرعت نسبی، اطلاعات موانع دینامیکی را تولید کردند. نشانگر هدف نیز با پردازش رنگ و نقشه اختلاف مکانی شناسایی شد و حسگرهای فراصوت، اطلاعات نزدیک‌برد مکمل را فراهم کردند. ارزیابی‌های شبیه‌سازی نشان دادند که دقت تخمین به فاصله، موقعیت جسم در میدان دید و حرکت ربات، به‌ویژه سرعت زاویه‌ای، وابسته است.

عامل یادگیری تقویتی بر پایه \lr{SAC} و مشاهده‌ای دیکشنری طراحی شد که از بردار ثابت ۱۵بعدی و ماتریس موانع \(10\times5\) تشکیل می‌شود. شبکه \lr{LSTM} نمایش فشرده موانع مرتب‌شده را استخراج و همراه ویژگی‌های ثابت به شاخه‌های کنشگر و دو منتقد منتقل می‌کند. فضای عمل پیوسته، فرمان‌های خطی و زاویه‌ای را تولید می‌کند و هموارسازی فرمان، تغییرات شدید را کاهش می‌دهد. تابع پاداش نیز پیشرفت به سمت هدف، زمان تا برخورد، هم‌راستایی تطبیقی، نرمی فرمان، فاصله فراصوتی، هزینه سپر و پیامدهای پایانی را ترکیب می‌کند. موفقیت تنها با قرارگیری در شعاع هدف، هم‌راستایی مناسب و نبود مانع در خط دید معتبر است.

سپر ایمنی، سطح خطر را از داده‌های فراصوت و موانع جلویی تعیین می‌کند. در پیکربندی آموزشی نهایی، سپر به‌صورت غیرفعال/نظارتی عمل می‌کند؛ بنابراین فرمان سیاست را بازنویسی نمی‌کند، بلکه سطح خطر را به مشاهده و هزینه‌ای با ضریب لاگرانژ تطبیقی تبدیل می‌نماید. در نهایت، آموزش معلم--دانشجو و برنامه مرحله‌ای، انتقال از هندسه دقیق به ادراک و سپس مواجهه با موانع و هدف متحرک را سامان دادند. بدین‌ترتیب، این فصل پیوند میان مبانی نظری و اجرایی سامانه را برقرار می‌کند و فصل بعد به پیاده‌سازی اختصاص دارد.
